
\documentclass[english,12pt]{article}
\usepackage[margin=1in]{geometry}
\usepackage{setspace}
\usepackage{amsmath}
\usepackage{amssymb}
\usepackage{amsthm}
\usepackage{float}
\usepackage[section]{placeins}
\usepackage{caption}
\DeclareCaptionFont{ninept}{\fontsize{9pt}{11pt}\selectfont}
\usepackage{afterpage}
\usepackage[T1]{fontenc}
\usepackage[utf8]{inputenc}
\usepackage{textcomp}
\usepackage{changepage}
\usepackage{color}
\usepackage{bbm}
\usepackage{graphicx}
\usepackage[dvipsnames]{xcolor}
\usepackage[colorlinks,citecolor=blue,linkcolor=black,urlcolor=black]{hyperref}
\usepackage{natbib}
\usepackage{appendix}
\usepackage{chngcntr}
\usepackage{booktabs}
\usepackage{subcaption}
\usepackage{pdflscape}
\usepackage[flushleft]{threeparttable}
\usepackage{rotating}
\newcommand{\ra}[1]{\renewcommand{\arraystretch}{#1}}
\usepackage{lmodern}
% CJE first submissions: main text double spaced; footnotes and references may be single spaced.
\usepackage[flushmargin]{footmisc}
\renewcommand{\footnotelayout}{\singlespacing}
\usepackage[none]{hyphenat}
\newtheorem{proposition}{Proposition}
\newtheorem{assumption}{Assumption}
\begin{document}
\sloppy
\title{When Demand Crosses Borders: Market Entry and the Cannabis Border Premium}
\author{Mahyar Rezazadeh\thanks{PhD Candidate, Department of Economics, Simon Fraser University. Email: \href{mailto:smr20@sfu.ca}{smr20@sfu.ca}}\\ {\small \emph{Simon Fraser University}}}
\date{\today}
\maketitle
\thispagestyle{empty}

\begin{abstract}
\noindent This paper asks how do local markets adjust to additional demand when firms are relatively free to enter? We study this question using recreational cannabis markets, where differences in legalization across U.S. states create spatially concentrated demand for legal retailers near borders with nonlegal states. Consumers in neighboring nonlegal states face lower travel costs when legal retailers are nearby, while firms decide whether and where to enter in response to the resulting market opportunity.

We first study cannabis market in New Mexico, where retailers entry is relatively unrestricted. The results show that ZIP codes bordering nonlegal-state lines have substantially more dispensaries and generate substantially more dispensary revenue, relative to population, than ZIP codes bordering states where recreational cannabis is already legal -- a comparison that isolates demand generated by captive cross-border access rather than generic proximity to any state line. Yet this market-level advantage does not translate into systematically higher revenue for individual dispensaries: store-level revenue outcomes are similar across the two groups. Taken together, these results are consistent with an equilibrium in which additional cross-border demand attracts entry and expands the size of the local market, rather than generating a persistent market power for incumbent retailers.

The composition of entry also differs across locations, but only briefly. A dynamic breakdown of the entry response shows that chain operators capture most of their eventual response to captive cross-border demand within the first quarter of legalization, while independent, single-location operators need roughly a year to catch up; after that, the two organizational forms respond identically. That transition window is exactly what our market-level entry game excludes by construction, since it is estimated on markets that have had time to settle. In that settled sample, the entry game -- built around two ex ante heterogeneous pools of potential entrants, independent and chain, choosing whether to enter each ZIP-code market, and estimated by simulated minimum distance -- finds that the independent pool's profitability responds more strongly to cross-border demand than the chain pool's does, at both the ZIP-code and, as a robustness check, the town level. We propose an own-network-internalization mechanism for this steady-state asymmetry -- an established operator's new outlet can divert sales from its own existing stores, while a first-location independent has no such network to internalize -- and test its central prediction directly using operator-level data. The direct test does not confirm the mechanism: the predicted interaction between an operator's existing network overlap and cross-border demand is not statistically distinguishable from zero. A separate dynamic exercise offers a candidate explanation for the early timing gap instead: independent and chain entrants convert opportunity into open stores at the same estimated speed once they act, but chains start from a much larger head start, consistent with existing financing and site-selection experience letting them act on day one rather than acting faster once underway. We report the null network-internalization result openly and discuss what it does, and does not, rule out.

Finally, we provide a suggestive comparison with Colorado, where local licensing rules constrain retail entry. Unlike New Mexico, Colorado border counties do not have more dispensaries than interior counties, while revenue per dispensary is higher near nonlegal-state borders. This contrast is consistent with the idea that restrictions on entry change how additional demand is absorbed: through market expansion when entry is relatively open, and through greater activity per incumbent when entry is constrained.

Overall, the findings show that the economic effects of spatially concentrated demand depend not only on the size of the demand opportunity, but on how easily firms can enter, and on which firms -- independent or chain -- choose to enter.

\medskip
\noindent \textbf{JEL classification:} L13, L51, H71, R12, K23 \\
\noindent \textbf{Keywords:} Cannabis legalization; Cross-border demand; Market entry regulation; Free entry; Strategic entry game; Organizational form; Border premium
\end{abstract}

\doublespacing
\newpage

\section{Introduction}


An increase in local demand can affect a market through more than one margin. Holding the number of firms fixed, greater demand may increase sales per incumbent. When entry is relatively open, however, higher expected returns can also attract new firms. The resulting entry expands market capacity and can limit the extent to which the initial demand advantage remains with individual retailers. This paper asks how local retail markets adjust to spatially concentrated demand when firms are relatively free to enter, using U.S. recreational cannabis markets where uneven legalization creates cross-border demand near nonlegal-state borders. Legalization has occurred unevenly across states, creating legal retail markets next to states where recreational sales remain prohibited. This generates a spatial source of demand such that consumers in nonlegal states who live close to a legal-state border face lower travel costs to reach a licensed retailer than consumers located farther away. A large literature shows that consumers respond to differences in prices, taxes, and regulation across jurisdictions and that these responses decline with travel costs \citep{AsplundFribergWilander2007,Lovenheim2008,ChandraHeadTappata2014}. In cannabis markets, \citet{HansenMillerWeber2020} show that sales at Washington retailers near Oregon fell sharply after Oregon opened its own recreational market, providing direct evidence that differences in legal access generate cross-border cannabis demand.

Our setting allows us to study a different margin of adjustment: whether cross-border demand is absorbed through additional firm entry rather than through higher revenue at incumbent retailers. \citet{HansenMillerWeber2020} examine incumbent retailers when cross-border demand disappears. We study market in New Mexico where firms can choose whether and where to enter while cross-border demand is present. Recreational sales began in 2022, and the state does not impose a statewide numerical cap on retail cannabis licenses. At the same time, large populations in neighboring Texas, Oklahoma, and Utah do not have access to recreational retail cannabis in their own states. Firms therefore make entry and location decisions in an environment in which cross-border demand is already present.

Consider a market in which consumers incur travel costs to reach a legal retailer and firms incur a fixed cost of entry. Proximity to consumers in a neighboring nonlegal state increases the effective demand faced by a border market. If the number of firms were fixed, this additional demand would increase revenue per retailer. When entry is relatively open, however, higher expected revenue also creates incentives for new firms to enter. As entry expands, the market can absorb the additional demand through a larger number of retailers, so aggregate market activity rises without necessarily generating a persistent revenue advantage for individual firms. This contrast between market-level expansion and firm-level outcomes motivates the empirical analysis.

We first compare New Mexico ZIP codes exposed to nearby populations in nonlegal states with ZIP codes exposed to nearby populations in states where recreational cannabis is already legal, Colorado and Arizona. Both groups of ZIP codes sit near a state line; they differ only in whether crossing that line gives a consumer legal access to cannabis at home. This comparison, developed further in Section~\ref{sec:framework}, isolates demand created by captive cross-border access from anything that is simply a feature of being close to a state boundary. Border exposure itself is defined using distance to the relevant state line and the population living on the far side. Our results show that ZIP codes near nonlegal-state lines generate substantially higher revenue and support substantially more dispensaries, relative to population, than ZIP codes near already-legal state lines. These differences are consistent with additional demand from consumers crossing into New Mexico to purchase cannabis legally, and, because both groups are equally close to a state line, are not easily explained by generic border proximity.

We then examine whether the larger market observed near nonlegal-state borders translates into better outcomes for individual retailers, but dispensaries in border locations do not generate systematically higher revenue than those near already-legal state lines, and revenue relative to the population surrounding each store is also similar across the two groups. Taken together, these findings suggest that the additional market activity near the border is absorbed through entry rather than through persistently higher revenue per retailer. Border markets therefore support more dispensaries and greater aggregate activity, while revenue per dispensary remains comparable to markets with no captive cross-border demand.

We next examine the composition of the entry response and find that it has a clear time signature: chain operators respond to captive cross-border demand almost immediately, independent, single-location operators respond with roughly a year's delay, and once both have had time to act, the two are statistically indistinguishable from each other. A simple diffusion model fit to the quarter-by-quarter response confirms that this is a difference in starting point rather than in speed -- both organizational forms close the remaining gap to their long-run response at essentially the same estimated rate -- consistent with a broader literature showing that multi-store retailers benefit from economies of scale and network economies that let them execute on a new opportunity quickly \citep{EllicksonHoughtonTimmins2013,HollenbeckGiroldo2022}, while independent and chain retailers otherwise differ in competitive position and location incentives \citep{AdamsEtAl2018,TalamasMarcos2025}. Restricting attention to markets that have had time to settle -- as our structural entry game does by construction -- we build a ZIP-code-level entry game with two ex ante heterogeneous pools of potential entrants, estimate the type-specific profitability response to cross-border demand in that settled steady state, and test a specific candidate explanation for why the two pools might differ there -- that established operators partly internalize the effect of a new outlet on their own nearby stores, while first-location independents do not. The direct operator-level test of this mechanism does not confirm it, and we say so plainly rather than treating the aggregate composition pattern as proof of a mechanism we cannot separately verify.

Finally, we provide a suggestive comparison with Colorado, where local governments have greater scope to restrict cannabis retail entry. Because the available Colorado data are less detailed and do not provide sufficiently precise measures of local regulatory constraints, we use this comparison as descriptive rather than causal evidence. Border counties in Colorado do not have systematically more dispensaries than interior counties, while revenue per dispensary is higher near nonlegal-state borders. This pattern is consistent with the entry framework, since a weaker response in the number of firms would leave a larger share of additional demand with existing retailers.

The paper contributes to four related literatures. First, it extends research on cross-border shopping by shifting attention from consumer responses and incumbent sales to the supply side of the market. While existing studies primarily examine how differences across jurisdictions affect where consumers purchase and how much incumbent firms sell, we show that persistent cross-border demand is also associated with changes in market size through additional entry.

Second, the paper contributes to research on entry and market structure. \citet{BresnahanReiss1991} and \citet{BerryReiss2007} emphasize that market size affects the number of firms a market can support, while entry in turn affects competitive outcomes. In cannabis markets, \citet{Thomas2019} and \citet{HollenbeckUetake2021} show that licensing restrictions can constrain entry and generate market power. We provide evidence from a relatively open-entry setting in which spatial differences in demand are reflected strongly in the number of firms rather than in revenue per individual retailer.

Third, we contribute to research on heterogeneity in retail entry and, specifically, to the literature on strategic store-network choice. \citet{Seim2006} and \citet{Mazzeo2002} model how potential entrants choose locations or product types anticipating rivals' choices; \citet{Nishida2008} and \citet{Holmes2011} model how a single organizational form's chain expands its network of outlets, trading off cannibalization of its own nearby stores against deterring rival entry. Our setting differs from this benchmark in two ways that shape the model in Section~\ref{sec:model}: markets face an observable, spatially discrete demand shifter -- captive cross-border demand -- that is largely absent from a local-demographics-driven setting, and two organizationally distinct pools of potential entrants compete for the same locations, so the question is not only how one chain balances cannibalization against deterrence, but why two different kinds of firms respond differently to the same demand shock in the first place.

Fourth, we contribute to research on heterogeneity in retail entry more broadly. The finding that chains and independents respond to captive cross-border demand at different speeds, converging to the same response within about a year, shows that the response to additional demand is not uniform across firm types even when it is eventually equal in size, complementing existing work on scale economies and organizational form by showing that spatially concentrated demand can affect the timing and composition of entry, not only its amount \citep{Sadun2015}.

Overall, the evidence suggests that cross-border demand affects the structure of the local cannabis market, not simply the revenue of retailers already operating there. In New Mexico, the additional demand associated with proximity to nonlegal states is reflected in greater aggregate activity and more dispensaries, while revenue per dispensary is similar across border and interior locations. The additional entry is also disproportionately associated with independent stores, a pattern the entry game in Section~\ref{sec:model} is built to explain and, in one specific dimension, fails to fully confirm. These findings highlight entry, and the composition of entry, as important margins through which local markets adjust to spatial differences in demand.

The rest of the paper is organized as follows. Section~\ref{litreview} reviews the related literature. Section~\ref{sec:framework} presents the conceptual framework and the two comparison groups used throughout. Section~\ref{sec:data} describes the data and institutional setting. Section~\ref{sec:empirical} presents the empirical strategy. Section~\ref{sec:result1} establishes that captive cross-border demand exists (Fact 1). Section~\ref{sec:mechanism} shows that markets respond primarily through entry rather than incumbent revenue (Fact 2). Section~\ref{sec:composition} shows that the entry response differs by organizational form (Fact 3) and motivates the structural model. Section~\ref{sec:model} builds the ZIP-code-level entry game. Section~\ref{sec:identification} discusses identification and estimation. Section~\ref{sec:modelfit} reports model fit. Section~\ref{sec:counterfactuals} reports counterfactuals and the mechanism test. Section~\ref{sec:falsification} presents robustness checks. Section~\ref{sec:discussion} discusses the Colorado comparison. Section~\ref{conclusion} concludes.

\section{Related Literature}
\label{litreview}

This paper relates most directly to four literatures: cross-border demand, market entry and regulation, strategic store-network choice, and heterogeneity across retail organizational forms.

The first literature studies consumer responses to differences across jurisdictions. Theoretical work by \citet{KanburKeen1993} shows how opportunities for cross-border shopping can affect competition between jurisdictions. Empirically, \citet{AsplundFribergWilander2007} find that alcohol sales respond to cross-border price differences, while \citet{Lovenheim2008} documents substantial cross-border cigarette purchases in response to interstate price differences. \citet{ChandraHeadTappata2014} show more generally that consumers weigh the gains from purchasing in another jurisdiction against the travel costs required to reach it. These studies establish a common result that is central to our setting: differences across jurisdictions generate spatially concentrated demand because the response is strongest where travel costs are lowest.

Related work shows that these demand differences can affect the location of retail activity. \citet{ManuszakMoul2009} find that gasoline tax differences alter the spatial distribution of retail capacity, indicating that cross-border shopping can affect not only where consumers buy but also where firms locate. In cannabis markets, \citet{HaoCowan2020} document spillovers from legalization into neighboring nonlegal states, while \citet{HansenMillerWeber2020} show directly that differences in recreational cannabis access generate cross-border retail demand. Our paper builds on this literature by examining how the number of firms responds when such demand is present during firms' entry and location decisions.

The second literature studies entry and market structure. \citet{BresnahanReiss1991} show that increases in market size support additional firms and that the competitive effects of entry can be large when markets initially contain few competitors. \citet{BerryReiss2007} review the broader empirical entry literature and emphasize the joint determination of profitability, entry, and market structure. The implication for our setting is that the effect of additional demand cannot be inferred from incumbent outcomes alone. If entry responds endogenously, a larger market can support more firms without generating substantially higher revenue for each firm.

Retail regulation can restrict this margin of adjustment. \citet{SchivardiViviano2011} show that barriers to retail entry affect market structure and incumbent outcomes. Cannabis markets provide particularly useful evidence because licensing rules directly constrain retail entry. \citet{Thomas2019} estimates a structural entry model for Washington's cannabis market and finds that license quotas substantially reduce entry relative to a free-entry counterfactual. \citet{HollenbeckUetake2021} show that these entry restrictions create local market power, while \citet{HollenbeckHristakevaUetake2026} find that additional nearby competitors reduce retail markups. Our contribution is complementary. Rather than studying a direct change in licensing rules, we examine how a relatively open-entry market responds to spatial differences in demand, and, going beyond this comparison group, build and estimate a discrete entry-threshold model of that response.

Third, we build directly on the strategic store-network-choice literature. \citet{Seim2006} and \citet{Mazzeo2002} develop entry-threshold models in which ex ante symmetric potential entrants within a type choose whether, and in Seim's case where or at what quality level, to enter, anticipating rivals' equilibrium behavior; both are estimated on a single cross-section of many isolated local markets, following the tradition of \citet{BresnahanReiss1991} and \citet{BresnahanReiss1990}. \citet{Nishida2008} and \citet{Holmes2011} instead model a single chain's network expansion decision directly, allowing a new outlet's profitability to depend on the chain's own existing nearby stores -- own-network cannibalization -- as well as on the presence of rival chains. Our model in Section~\ref{sec:model} combines elements of both traditions: an entry-threshold game between two ex ante heterogeneous pools of potential entrants, in the spirit of Seim and Mazzeo, in which one pool's type-specific profitability is derived from an own-network-internalization mechanism in the spirit of Nishida and Holmes. Unlike Nishida's or Holmes's settings, our markets also carry an observable, spatially discrete demand shifter -- captive cross-border demand -- that lets us test whether the two pools' equilibrium entry differs in a way a demand-side story, rather than a network-choice story alone, would predict.

The fourth literature concerns heterogeneity across potential entrants more broadly. \citet{EllicksonHoughtonTimmins2013} show that retail chains benefit from network economies that affect location decisions. In cannabis retailing, \citet{HollenbeckGiroldo2022} find substantial returns to operating multiple stores, indicating meaningful economies of scale among multi-store firms. At the same time, organizational forms can differ in their local competitive advantages. \citet{AdamsEtAl2018} find that independent and chain coffee retailers occupy different competitive positions, while \citet{TalamasMarcos2025} documents characteristics that can allow owner-operated neighborhood retailers to compete with larger chains. \citet{Sadun2015} further shows that regulation can change not only the amount of retail entry but also the form through which firms enter.

Our finding that chain and independent cannabis retailers respond to captive cross-border demand on different timelines, and, once the market has had time to settle, at different steady-state rates, adds a spatial demand dimension to this literature, and Section~\ref{sec:model} takes the further step of proposing, and directly testing, a specific mechanism for the steady-state difference. The direct test does not confirm the proposed mechanism, so the paper does not claim to have identified why the entry response differs across organizational forms; it establishes that the response differs both in timing and, once settled, in level, derives one economically motivated candidate explanation for the latter, and reports honestly that this explanation is not the one the data confirm. This distinction matters because a market's response to additional demand can involve changes in the timing, number, and composition of firms all at once, and because a coherent candidate mechanism is not the same thing as an identified one.

Taken together, these literatures motivate our empirical approach by linking four elements of the analysis: differences in legal access create cross-border demand, travel costs concentrate that demand near state borders, endogenous entry determines whether the resulting market opportunity appears through additional firms or through higher revenue for existing retailers, and, where entry responds, an explicit entry-threshold model of heterogeneous potential entrants is needed to say whether and why the response differs by organizational form. New Mexico provides a setting in which this full sequence can be observed under relatively open entry, while the Colorado comparison offers suggestive evidence on how the same demand may be absorbed differently when entry is more constrained.

\section{Conceptual Framework}
\label{sec:framework}

Before turning to the data, it helps to set out what a simple model of store entry predicts about cross-border demand, and to state precisely which two groups of ZIP codes that prediction is tested on. The model has three ingredients. Consumers are spread across space and pay a travel cost to reach a store, each store pays a fixed cost to operate, and firms keep entering as long as they expect to cover that fixed cost. This is the free-entry spatial setup of \citet{Salop1979}, applied to the entry question studied by \citet{BresnahanReiss1991}.

\paragraph{Setup.} A local market is served by stores spread along a line, with consumers at a roughly constant density. Each consumer buys one unit from the nearest store and pays the posted price plus a travel cost that grows with distance. Let $S$ be the size of local demand, meaning the number of customers the market can serve in a period. A market's demand $S$ is larger when it can draw on consumers who live across a state line where recreational cannabis is not legally available, since those consumers would otherwise have no legal in-state option. A market near a state line where cannabis is \emph{already} legal draws on no such captive population: consumers on the far side already have a legal retailer of their own. Writing $NL$ (nonlegal-border) and $LB$ (legal-border) for the two groups of ZIP codes compared throughout the empirical sections, $S_{NL} > S_{LB}$, while a market with no state-line exposure at all, an interior ZIP code, has $S_I \lessapprox S_{LB}$: both groups are equally close to \emph{some} state line, but only $NL$ carries the extra captive demand. Firms first decide whether to enter and where to locate, then compete on price.

\paragraph{Free entry.} When several stores compete, each store's price settles at a markup over cost that is governed by how far apart neighboring stores are, that is, by the local density of stores, and not by the overall size of the market. As a result, the store size at which a store just breaks even, call it $\bar q$, and the break-even price, call it $\bar p$, are pinned down by the travel cost, the fixed cost, and the local consumer density, but not by $S$.%
\footnote{In the circular version of the model with consumers at density $\rho$ per unit of distance, the symmetric free-entry equilibrium has $\bar q = \sqrt{\rho f / t}$ and $\bar p = c + \sqrt{t f / \rho}$, where $t$ is the travel cost per unit distance, $f$ the per-store fixed cost, and $c$ the marginal cost of serving a customer. Neither expression depends on $S$ once a larger market is read as a larger catchment area at similar density rather than a denser population.}
Entry then continues until the typical store just breaks even:
\begin{equation}
\label{eq:zeroprofit}
(\bar p - c)\,\bar q = f .
\end{equation}
Because $\bar q$ does not depend on $S$, the number of stores the market supports is
\begin{equation}
\label{eq:nstar}
n^\ast = \frac{S}{\bar q},
\end{equation}
which rises one for one with local demand. Total market revenue is $\bar p\, S$, so it also rises with demand and, expressed per resident, is higher wherever demand per resident is higher. Revenue per store is $\bar p\, \bar q$, which does not move with $S$.

Put plainly, when entry is free the extra demand at a nonlegal-border location is met by adding stores, not by making the existing stores bigger or letting them charge more. Captive cross-border demand shows up as a denser and larger market rather than as fatter margins for incumbents.\footnote{Letting stores set prices strategically rather than charge a fixed break-even markup only reinforces this. With more stores nearby, each store faces more elastic demand and its markup falls, so entry both splits demand among more sellers and pushes margins down. Both forces work against a rise in revenue per store when captive demand pulls in new entrants, so P2 is a conservative test of the mechanism.}

\paragraph{Constrained entry.} Now suppose a regulator caps the number of stores at $\bar n$, below the free-entry level that captive demand would otherwise support. The cap does not change how much consumers want to buy, so each of the $\bar n$ incumbents serves more customers than the break-even level, $S_{NL} / \bar n > \bar q$. With fewer competitors, price is not competed down as far, so revenue per store rises above $\bar p\,\bar q$ and incumbents earn a rent. Dispensary density still rises with captive demand, but by less than one for one, and part of the extra demand is captured as higher revenue per store rather than as new entry. The more tightly the cap binds, the smaller the density response and the larger the incumbent rent.

The gap between the free-entry density response and the constrained one is informative on its own. It measures how many stores the entry barrier keeps out of exposed markets. Multiplying that shortfall by revenue per store converts it into the legal-market revenue, and the associated tax base, that the barrier leaves on the table.

\paragraph{Why entrants may differ, and what the model has to say about it.} The free-entry logic above is silent on \emph{who} the entrants are. Section~\ref{sec:composition} documents that they are not drawn uniformly from the pool of potential firms: independent, first-location entrants respond to captive cross-border demand more strongly than established multi-store operators do. A market-size argument alone does not explain this, since a bigger market attracts more of \emph{any} type of firm. Section~\ref{sec:model} extends the entry-threshold logic of this section to two ex ante heterogeneous pools of potential entrants and derives a primitive -- an established operator's own-network overlap with the candidate market -- under which the two pools' profitability responds differently to the same demand shock even though they face the same underlying market-size expansion described here.

\paragraph{Predictions.} The framework gives four predictions that map onto the empirical sections.

\begin{description}
\item[P1 (market level, free entry).] Where entry is free, captive cross-border demand -- proximity to a nonlegal-state line rather than to any state line -- raises dispensary density and total cannabis revenue per resident.
\item[P2 (firm level, free entry).] Where entry is free, captive cross-border demand does not raise revenue per store. The additional demand is absorbed by more stores, not by larger or higher-priced ones.
\item[P3 (composition).] If independent and chain potential entrants differ in how much of a new outlet's revenue is incremental to the firm, as Section~\ref{sec:model} formalizes, the entry response to captive demand need not be uniform across organizational forms.
\item[P4 (constrained entry).] Where entry is capped below the free-entry level, the density response to captive demand is weaker and part of the extra demand appears as higher revenue per incumbent store. The difference between the free-entry density response in New Mexico and the constrained response in Colorado measures how much the entry barrier binds, and applied to revenue per store it gives the legal-market revenue and tax base forgone.
\end{description}

Table~\ref{tab:predictions} shows where each prediction is tested. P1 and P2 carry the main analysis and are tested on New Mexico, comparing nonlegal-border ZIP codes to always-legal-border ZIP codes as the primary treatment and control groups; a pooled control that also includes interior ZIP codes, and a town-level aggregation, are reported as robustness checks. P3 motivates and is tested by the structural model of Section~\ref{sec:model}. P4 is a more tentative extension, because the cross-state comparison behind it uses only two states and yields a marginally significant estimate.

\begin{table}[htbp]
\centering
\caption{Predictions and where they are tested}
\label{tab:predictions}
\begin{tabular}{@{}p{0.07\textwidth} p{0.55\textwidth} p{0.30\textwidth}@{}}
\toprule
& Empirical counterpart & Where \\
\midrule
P1 & Nonlegal-border coefficient on dispensary count and revenue in the New Mexico ZIP-code panel (equation~\ref{baseline}) & Section~\ref{sec:result1} \\[4pt]
P2 & Nonlegal-border coefficient on revenue per store, ZIP-level decomposition and dispensary-level panel & Section~\ref{sec:mechanism} \\[4pt]
P3 & Nonlegal-border coefficient on independent and chain entrant counts separately (equation~\ref{eq:composition}); estimated type-specific demand response $\hat\delta_I$ vs.\ $\hat\delta_C$ in the entry game & Section~\ref{sec:composition}, Section~\ref{sec:modelfit} \\[4pt]
P4 & $\text{Border}\times\text{Restricted}$ interaction in the New Mexico versus Colorado county panel, and the implied foregone-revenue counterfactual & Section~\ref{sec:discussion}, Appendix~\ref{sec:crossstate} \\
\bottomrule
\end{tabular}
\end{table}

\paragraph{Taking the model to the data.} Three points are worth stating up front. First, the model speaks to revenue rather than to prices directly, because the data contain dispensary revenue but not transaction prices. P2 is therefore a test of whether revenue per store rises with captive demand, read as the combined effect of store size and price. Second, the model treats a captive-demand market as a larger catchment area at comparable consumer density rather than as a denser population; this matches the setting, where nonlegal- and legal-border ZIP codes have broadly comparable resident populations. Third, and specific to this revision: the main comparison throughout is nonlegal-border versus always-legal-border ZIP codes, not border versus interior. Both groups of ZIP codes sit near a state line; only the nonlegal-border group has access to a captive population with no legal in-state alternative. This comparison is a sharper test of P1--P3 than a border-versus-interior comparison, because it holds fixed whatever is generic about operating a retail business near \emph{any} state boundary -- transportation access, land values, or unmeasured local commercial activity -- and isolates the piece that is specific to prohibition across the line.

\section{Data and Background}
\label{sec:data}

The empirical setting combines two legal cannabis markets with different approaches to retail entry. New Mexico is the primary setting for the analysis. Recreational cannabis sales began in April 2022, and state law does not impose a numerical cap on retail licenses. Licensing is largely administrative, allowing the number and location of dispensaries to respond to differences in local demand. Colorado provides a useful comparison. Although Colorado does not impose a statewide numerical cap, local governments may restrict or prohibit cannabis businesses, making entry more constrained in many parts of the state. This institutional contrast makes it possible first to examine how markets exposed to captive cross-border demand develop when entry is relatively open and then to assess whether the same comparison looks different where entry is subject to greater local restrictions.

The primary unit of analysis is the ZIP code, consistent with the design used throughout this paper. \textbf{The main empirical analysis uses the full ZIP-code-by-quarter panel for New Mexico: all 368 populated ZIP codes from 2022Q2 through 2026Q2, with true zero dispensary activity retained rather than dropped.} A ZIP code that never records a dispensary is an economically meaningful outcome, not a missing observation, and restricting to ZIP codes that happen to have activity would mechanically exclude the extensive-margin entry response -- whether a market gets a dispensary at all -- that is the central object of this paper's free-entry mechanism. As a secondary check, the paper also reports the same specifications restricted to the 125 ZIP codes that record at least one operating dispensary at some point in the sample; this conditional-on-activity sample speaks to the intensive margin among markets that already have some activity, and is reported alongside every main result rather than in place of it. To study outcomes at the individual-retailer level, the paper also uses a dispensary-by-quarter panel of geocoded New Mexico dispensaries over the same period. Town-level and county-level aggregations of the same underlying storefront data are used only as robustness checks, reported alongside the corresponding ZIP-code result: the town aggregation is itself a full zero-inclusive panel of all New Mexico town-markets (excluding the Albuquerque metro, reported separately) following this project's own structural-modeling convention, and the cross-state comparison uses county-level data from New Mexico and Colorado (Section~\ref{sec:crossstate}).

The New Mexico ZIP and dispensary panels are balanced. Because sales are reported monthly at the licensee level, every populated ZIP code is observed in every quarter, allowing true zero activity to be distinguished from missing observations. This distinction is important for the main outcomes because a ZIP code with no operating dispensary represents an economically meaningful outcome rather than an observation that should be excluded; the main specifications reported throughout Sections~\ref{sec:result1}--\ref{sec:composition} retain these ZIP-quarters directly using the inverse hyperbolic sine transformation, and the conditional-on-activity secondary check (Section~\ref{sec:falsification}) is reported alongside them. The county panel is less complete: counties enter the sample once dispensary activity is observed and may disappear in periods in which revenue is suppressed for confidentiality.

Each dispensary's ZIP code is its own reported location, matched to the 368-ZIP crosswalk directly; a small number of storefronts (13 of 1,255) carried a self-reported ZIP code that does not correspond to any real New Mexico ZIP -- a transposed digit, a P.O.-box code, or, in one case, a wrong-state prefix -- and each of these, and only these, is assigned to its nearest real ZIP by geographic distance. Every other storefront keeps its own reported ZIP code unchanged.

Border exposure is defined using geography rather than observed dispensary locations, so that a market's exposure status is determined independently of firms' entry decisions. A ZIP code, dispensary, or county is classified as exposed to a given state line when two conditions are satisfied. First, its centroid, or the dispensary's own location, must lie within 75 miles of the relevant state line. Second, at least 1,000 residents of the neighboring state must live within the same 75-mile area. The population requirement excludes sparsely populated sections of the state line for which physical proximity does not imply economically meaningful access to potential cross-line consumers.

Two lines matter for New Mexico. \textbf{Nonlegal-border} ZIP codes satisfy the 75-mile/1,000-resident rule on the Texas, Oklahoma, or Utah line: consumers on the far side of these lines have no legal recreational retailer at home, so a nonlegal-border New Mexico dispensary can capture genuinely captive demand. \textbf{Always-legal-border} ZIP codes satisfy the same rule on the Colorado or Arizona line, both of which operated recreational retail markets throughout the sample period: consumers on the far side already have a legal option at home, so an always-legal-border New Mexico dispensary faces no comparable captive pool, even though it is, by construction, exactly as close to a state line as a nonlegal-border dispensary is. A ZIP code satisfying both rules is classified as nonlegal-border: 14 of New Mexico's 368 populated ZIP codes, mostly near the Four Corners area and the state's northeastern corner, lie within 75 miles of both a nonlegal- and an already-legal-state line, and the nonlegal rule takes priority for these. The remaining, non-exposed ZIP codes are \textbf{interior}. Under this classification, 92 of New Mexico's 368 populated ZIP codes are nonlegal-border, 117 are always-legal-border, and 159 are interior; among the 125 ZIP codes with any recorded dispensary activity, 35 are nonlegal-border, 24 are always-legal-border, and 66 are interior.

\textbf{The main comparison throughout this paper is nonlegal-border versus always-legal-border ZIP codes, estimated on the full 368-ZIP-code panel}, motivated in Section~\ref{sec:framework}: both groups sit near a state line, so the comparison isolates the demand effect of captive, prohibition-driven cross-border access from anything that is simply a feature of border proximity, and the full panel keeps every ZIP-quarter in which a market never attracts a dispensary rather than dropping it. A pooled control group -- always-legal-border ZIP codes together with interior ZIP codes -- a full zero-inclusive town-level aggregation, and the conditional-on-activity secondary sample (125 ZIP codes with recorded dispensary activity) are reported alongside every main result, and Albuquerque is included or excluded as noted in each case, since it is interior by construction and can only affect specifications that use the pooled control.

Figure~\ref{fig:spatial_inequality} compares dispensary density and revenue per resident across ZIP codes in New Mexico; Appendix Figure~\ref{fig:setting} places New Mexico alongside neighboring Colorado.

\begin{figure}[htbp]
\centering
\includegraphics[width=\textwidth]{fig_nm_spatial_inequality.png}
\caption{Spatial inequality in dispensary access, New Mexico. Panel (a): each
bubble represents a ZIP-code centroid, sized by average dispensaries per 1,000
residents over 2022Q2--2026Q2 and colored by border status. Panel (b):
population-weighted dispensaries per 1,000 residents (solid bars, left axis) and
revenue per resident (hatched bars, right axis) for border and interior ZIP
codes, pooling all ZIP-quarters.}
\label{fig:spatial_inequality}
\end{figure}

New Mexico cannabis sales are obtained from the Cannabis Control Division's Cannabis Reporting Online Portal \citep{NMCCDReporting}, which reports monthly gross sales by individual licensee and municipality. These data are aggregated to the dispensary, ZIP-code, town, and county levels according to the unit of analysis. Colorado sales are obtained from the Colorado Department of Revenue's county-level marijuana sales series \citep{COMarijuanaSalesByCounty}. Dispensary locations in New Mexico are obtained from the Cannabis Control Division's public list of licensed cannabis establishments \citep{NMCCDLicensees}, geocoded using the U.S. Census Bureau's batch geocoder. Colorado dispensary information is obtained from the Marijuana Enforcement Division's licensee records \citep{COMedLicensees}. Each dispensary's organizational type -- independent or chain -- is determined from the same licensee records: an operator (brand) that has ever run two or more storefronts in New Mexico over the sample is classified as a chain; an operator that has run only one is classified as independent. Section~\ref{sec:composition} classifies each store's type \emph{at the time it enters}, using only information available at that storefront's own first active quarter, so that an independent's later expansion into a chain does not retroactively change how its original entry is coded.

Population and geographic data are obtained from the SimpleMaps U.S. ZIP database \citep{SimpleMapsZIP}, the U.S. Census Bureau's Population Estimates Program \citep{CensusPopEst2023}, 2020 Census block populations \citep{CensusBlocks2020}, and the Census Bureau's cartographic boundary files \citep{CensusCountyShapefile}. For the falsification exercises in Section~\ref{sec:falsification}, the paper also uses non-cannabis retail data from the New Mexico Taxation and Revenue Department \citep{NMTRDGrossReceipts} and the Census Bureau's ZIP Codes Business Patterns \citep{CensusZBP}.

The analysis focuses on four main outcomes. Revenue per resident and dispensary density (per 1,000 residents) characterize the overall size of the local cannabis market and are computed on the full, zero-inclusive ZIP panel using the inverse hyperbolic sine transformation, since a ZIP code with no dispensary is a valid zero rather than a missing value. Revenue per dispensary measures the sales generated by the average retailer and is defined only for markets containing at least one store, so it is unavoidably conditional on activity regardless of which sample is used. The number of independent and the number of chain dispensaries, defined at entry, characterize the composition of the market and are likewise reported on the full zero-inclusive panel. When outcomes enter the regressions in logarithms (the conditional-on-activity secondary specifications only), coefficients are converted into percentage differences using $(e^{\hat\beta}-1)\times100$. Appendix Table~\ref{tab:sumstats_nm} summarizes the main New Mexico variables across the full 368-ZIP-code sample and the three exposure groups, alongside the same summary restricted to the 125 ZIP codes with recorded dispensary activity, and Appendix Table~\ref{tab:sumstats_state} summarizes the cross-state county sample.

The data have several limitations. First, the New Mexico panels begin when recreational sales commenced in April 2022, so the analysis compares exposed and unexposed markets after legalization rather than estimating the causal effect of legalization itself. Second, population is measured using a fixed vintage and is not adjusted for the legal purchase age. Third, Colorado's historical dispensary counts are reconstructed from license records rather than observed in a complete operating panel. Fourth, the data do not identify the residence of individual cannabis consumers; cross-border demand is measured using the size and proximity of populations on the far side of the relevant state line rather than observed directly. Finally, nonlegal- and always-legal-border ZIP codes can differ in when their first dispensary opened -- always-legal-border ZIP codes are, on average, younger markets -- so the conditional-on-activity secondary specification controls for market age (quarters since each ZIP code's own first active dispensary) in addition to calendar-quarter fixed effects. The main, full-panel specification does not include a market-age control, since market age is undefined for a ZIP code that never has a dispensary; a version of the market-age control that buckets never-active ZIP-quarters separately was checked and rejected because it absorbs, rather than measures, the extensive-margin entry response that is the object of this paper's main analysis, a point documented further in Section~\ref{sec:empirical}.

\section{Empirical Strategy}
\label{sec:empirical}

The estimand throughout Sections~\ref{sec:result1}--\ref{sec:composition} is the effect of captive cross-border demand -- proximity to a state line where recreational cannabis is not legally available, relative to proximity to a state line where it already is -- on outcomes in New Mexico's dispensary market. Legalization created the legal recreational market at a discrete point in time; because the timing was not fully anticipated by firms or consumers, the policy generated a relatively sharp market-opening shock, after which nonlegal-border ZIP codes were able to serve an additional pool of consumers that always-legal-border ZIP codes could not.

\subsection{Border-Exposure Panel Design}
\label{sec:borderexposure}

The main specification compares outcomes across ZIP codes that satisfy the 75-mile/1,000-resident rule on a nonlegal-state line against ZIP codes that satisfy the same rule on an already-legal-state line, estimated on the full, zero-inclusive 368-ZIP-code panel:

\begin{equation}
\label{baseline}
\operatorname{asinh}(y_{zt})=\alpha
+
\beta\, \text{Nonlegal}_z
+
\gamma \log(\text{Population}_z)
+
\delta_t
+
\varepsilon_{zt}.
\end{equation}

The outcome $y_{zt}$ is one of the market-level cannabis outcomes in ZIP code $z$ and quarter $t$: revenue per resident, dispensary density (per 1,000 residents), or, in Section~\ref{sec:composition}, the number of independent and the number of chain dispensaries. The inverse hyperbolic sine transformation is used in place of the log because a ZIP-quarter can genuinely record zero dispensaries and zero revenue, and dropping those observations would mechanically remove the extensive-margin variation -- whether a market gets a dispensary at all -- that the free-entry mechanism (Section~\ref{sec:framework}) predicts responds to captive demand. $\text{Nonlegal}_z$ equals one if $z$ is a nonlegal-border ZIP code and the sample is restricted to nonlegal- and always-legal-border ZIP codes only, so $\beta$ is the coefficient of interest -- the average difference in the outcome attributable to captive cross-border demand, holding fixed proximity to a state line. $\text{Population}_z$ is the resident population of ZIP code $z$; its coefficient $\gamma$ captures how the outcome scales with the local, legal-side market and is not itself a measure of cross-border demand. $\delta_t$ denotes calendar-quarter fixed effects, absorbing statewide seasonality and market growth. Standard errors are clustered at the county level. Revenue per dispensary, unavoidably defined only where a ZIP-quarter has at least one operating store, is estimated in logs on the same full-panel ZIP-quarters that have positive store counts.

Equation~\eqref{baseline} is re-estimated under three additional samples, reported alongside the main result in every table. The \textbf{pooled-control} sample restores the interior ZIP codes dropped from the main specification, redefining $\text{Nonlegal}_z$ against always-legal-border and interior ZIP codes combined; this is the design closer to a conventional border-versus-interior comparison and is reported because it is, in most outcomes, a larger and better-powered sample, at some cost to the sharper captive-demand interpretation of the main specification. The \textbf{town-level} aggregation re-estimates equation~\eqref{baseline} on a full zero-inclusive panel of New Mexico town-markets, with Albuquerque excluded by the same convention used throughout this project's structural modeling (Section~\ref{sec:model}). The \textbf{conditional-on-activity} secondary specification restricts the sample to the 125 ZIP codes with at least one recorded dispensary, replaces $\operatorname{asinh}(y_{zt})$ with $\log(y_{zt})$ (valid here since these ZIP-quarters are conditioned to be positive for the store-count and revenue outcomes), and adds a market-age control $\theta_{\text{age}(z,t)}$ -- one indicator for each value of quarters elapsed since ZIP code $z$'s own first active dispensary -- because, conditional on already having a dispensary, always-legal-border ZIP codes are on average younger markets than nonlegal-border ones, so that calendar-quarter fixed effects alone would partly confound market vintage with captive-demand exposure within this conditional sample. The market-age control is not used in the main, full-panel specification: market age is undefined for a ZIP code that never has a dispensary, and bucketing never-active ZIP-quarters into their own indicator was checked and found to absorb, rather than control for, the extensive-margin entry response itself (Section~\ref{sec:falsification}).

Composition equation~\eqref{eq:composition} is equation~\eqref{baseline} applied separately to the independent and chain dispensary counts, with the same $\operatorname{asinh}$ transformation in the main specification (since a ZIP-quarter can record zero independent or zero chain dispensaries) and the same log-plus-market-age alternative in the conditional-on-activity secondary specification:
\begin{equation}
\label{eq:composition}
\operatorname{asinh}(N^\tau_{zt})=\alpha_\tau
+
\beta_\tau\, \text{Nonlegal}_z
+
\gamma_\tau \log(\text{Population}_z)
+
\delta_{\tau,t}
+
\varepsilon_{\tau,zt}, \qquad \tau \in \{I, C\},
\end{equation}
estimated as two separate regressions rather than a single differenced outcome, so that $\hat\beta_I$ and $\hat\beta_C$ can each be read directly and compared to the entry game's own $\hat\delta_I$ and $\hat\delta_C$ in Section~\ref{sec:modelfit}.

The identifying assumption is that, absent captive cross-border demand, nonlegal- and always-legal-border ZIP codes would exhibit similar cannabis-market outcomes. This assumption is not automatic: the two groups could differ in remoteness, local commercial activity, or other characteristics correlated with which international or interstate line they happen to sit near. Section~\ref{sec:falsification} reports several checks -- a continuous distance/population measure in place of the discrete $\text{Nonlegal}_z$ indicator, 50- and 100-mile alternative thresholds, non-cannabis placebo sectors, and the pooled-control, town-level, and conditional-on-activity samples already introduced above.

\section{Fact 1: Captive Cross-Border Demand Exists}
\label{sec:result1}

Table~\ref{tab:fact1fact2} reports the baseline estimates of equation~\eqref{baseline} for revenue per resident, dispensary density, and revenue per store, on the full 368-ZIP-code panel. In the main specification -- nonlegal-border versus always-legal-border ZIP codes, 3,553 ZIP-quarters -- revenue per resident and dispensary density are both significantly higher in nonlegal-border ZIP codes ($\hat\beta=+0.835$, $p<0.01$, and $\hat\beta=+0.109$, $p<0.05$, respectively, in inverse-hyperbolic-sine units). Widening the control group to always-legal-border and interior ZIP codes pooled together (6,256 ZIP-quarters) gives the same pattern, somewhat smaller in magnitude but still significant at conventional levels ($\hat\beta=+0.575$, $p<0.05$, for revenue per resident; $\hat\beta=+0.096$, $p<0.05$, for density).

The secondary conditional-on-activity specification -- restricting to the 125 ZIP codes with recorded dispensary activity, in logs, with a market-age control -- gives the same qualitative pattern using un-normalized outcomes (Appendix Table~\ref{tab:fact1fact2conditional}): total dispensary revenue is 89.0 percent higher $\left(e^{0.637}-1\right)$ and the number of operating dispensaries is 57.0 percent higher $\left(e^{0.451}-1\right)$ in nonlegal-border ZIP codes, both statistically distinguishable from zero. The full-panel main result and the conditional-on-activity secondary result therefore agree on the qualitative finding; the full panel is preferred because, unlike the conditional sample, it can also detect the extensive-margin difference in whether a ZIP code ever attracts a dispensary at all, rather than only the intensive-margin difference among ZIP codes that already have one.

Because both nonlegal- and always-legal-border ZIP codes are, by construction, within 75 miles of some state line, this comparison is not explained by generic proximity to a border -- transportation access, land values, or other characteristics of border geography that have nothing to do with cannabis prohibition across the line. The difference is therefore read as evidence of demand generated specifically by access to consumers for whom recreational cannabis remains illegal at home, consistent with P1.

\begin{table}[htbp]
\centering\scriptsize
\begin{tabular}{@{}lcccccc@{}}
\toprule
Quarter & \multicolumn{2}{c}{Revenue per resident (asinh)} & \multicolumn{2}{c}{Dispensary density (asinh)} & \multicolumn{2}{c}{log(Revenue per store)} \\
\cmidrule(lr){2-3} \cmidrule(lr){4-5} \cmidrule(lr){6-7}
& Coef. & (SE) & Coef. & (SE) & Coef. & (SE) \\
\midrule
2022Q2 & +0.459** & (0.203) & +0.017 & (0.034) & -0.355 & (0.753) \\
2022Q3 & +0.583** & (0.240) & +0.039 & (0.029) & -0.322 & (0.238) \\
2022Q4 & +0.729*** & (0.218) & +0.063** & (0.030) & +0.094 & (0.331) \\
2023Q1 & +0.832*** & (0.252) & +0.084** & (0.035) & +0.298 & (0.370) \\
2023Q2 & +0.868*** & (0.257) & +0.105*** & (0.036) & +0.231 & (0.287) \\
2023Q3 & +0.904*** & (0.280) & +0.117** & (0.048) & +0.415 & (0.414) \\
2023Q4 & +0.944*** & (0.292) & +0.135*** & (0.049) & +0.348 & (0.328) \\
2024Q1 & +0.951*** & (0.295) & +0.152*** & (0.058) & +0.286 & (0.347) \\
2024Q2 & +0.954*** & (0.304) & +0.166*** & (0.058) & +0.337 & (0.370) \\
2024Q3 & +0.961*** & (0.306) & +0.171*** & (0.058) & +0.332 & (0.506) \\
2024Q4 & +0.960*** & (0.293) & +0.174*** & (0.055) & +0.229 & (0.192) \\
2025Q1 & +0.900*** & (0.308) & +0.145** & (0.070) & +0.519 & (0.341) \\
2025Q2 & +0.818*** & (0.303) & +0.114* & (0.068) & +0.336 & (0.214) \\
2025Q3 & +0.818*** & (0.312) & +0.114* & (0.063) & +0.027 & (0.239) \\
2025Q4 & +0.846*** & (0.304) & +0.097 & (0.061) & +0.129 & (0.234) \\
2026Q1 & +0.817*** & (0.309) & +0.085 & (0.062) & +0.186 & (0.229) \\
2026Q2 & +0.858*** & (0.299) & +0.081 & (0.062) & +0.664 & (0.442) \\
\bottomrule
\end{tabular}

\caption{Fact 1 and Fact 2, dynamic specification: quarter-specific nonlegal- vs.\ always-legal-border coefficients $\hat\beta_q$ for revenue per resident, dispensary density, and revenue per store. Sample: full New Mexico ZIP-code panel, 2022Q2--2026Q2 (368 ZIP codes, zero-inclusive; revenue-per-store column restricted to ZIP-quarters with at least one operating dispensary). Each $\hat\beta_q$ is the nonlegal-vs.\ always-legal-border gap in that quarter, not a deviation from an omitted reference quarter. Revenue-per-resident and density columns in inverse-hyperbolic-sine units; revenue-per-store column in logs. All quarters include $\log(\text{population})$ and calendar-quarter fixed effects; standard errors (in parentheses) clustered by county. $^{*}p<.10$ $^{**}p<.05$ $^{***}p<.01$.}
\label{tab:fact1fact2}
\end{table}

\FloatBarrier

\section{Fact 2: Markets Respond Through Entry}
\label{sec:mechanism}

Per-capita revenue is the product of dispensaries per capita and revenue per dispensary, so the ZIP-level result in Table~\ref{tab:fact1fact2} can be decomposed exactly into an extensive margin (dispensary density) and an intensive margin (revenue per store). Revenue per store is 29.7 percent higher in nonlegal- relative to always-legal-border ZIP codes in the main specification, but the estimate is not statistically distinguishable from zero. In the pooled-control sample it is larger, 32.8 percent, and statistically significant at the 5 percent level. Comparing the two margins, density is precisely estimated in both samples while revenue per store is either insignificant or only marginally so: the extensive margin, not the intensive one, is the more reliable channel behind the ZIP-level revenue difference documented in Fact 1, consistent with P2.

This ZIP-level decomposition is corroborated by the dispensary-level panel, which follows individual stores rather than ZIP-code aggregates and controls for brand fixed effects. Figure~\ref{fig_dispensary} plots the dynamic dispensary-level border coefficient on revenue per local resident under the legal-border control group; the estimates fluctuate around zero with no systematic upward pattern, indicating that captive cross-border demand does not translate into persistently higher revenue for individual incumbent stores. Figure~\ref{fig:dynamic_legalcontrol} shows the corresponding ZIP-level dynamic pattern under the nonlegal- versus always-legal-border comparison: the gap in dispensary density and revenue per capita is large and persistent across most of the sample period, while the gap in revenue per dispensary is smaller and inconsistently significant, the same pattern as the static estimates in Table~\ref{tab:fact1fact2}.

\begin{figure}[!ht]
\centering
\includegraphics[width=\textwidth]{fig_nm_dynamic_legalcontrol.png}
\captionsetup{font=ninept}
\caption{Dynamic captive cross-border demand estimates, nonlegal-border ZIP codes versus the always-legal-border control group. Nonlegal-border ZIP codes: 38; always-legal-border ZIP codes: 19 (116-ZIP sample vintage). County-clustered 95\% confidence intervals, reference quarter 2022Q2.}
\label{fig:dynamic_legalcontrol}
\end{figure}

\begin{figure}[!ht]
\centering
\includegraphics[width=\textwidth]{fig_eventstudy_nm_dispensary_revperresident_1km_vs_2km.png}
\captionsetup{font=ninept}
\caption{Revenue per local resident, dispensary-level panel, one- versus two-kilometre denominator. County-clustered 95 percent confidence intervals.}
\label{fig_dispensary}
\end{figure}

Together, the ZIP-level decomposition and the dispensary-level panel support the same conclusion from two different levels of aggregation: captive cross-border demand is absorbed primarily by an increase in the number of operating dispensaries, not by a persistent increase in revenue at the stores that are already there. This is the pattern predicted by a free-entry model (Section~\ref{sec:framework}): additional demand attracts entrants until the typical store is again near its break-even size, so the demand advantage shows up as market breadth rather than incumbent rents.

\FloatBarrier

\section{Fact 3: Entrant Composition Responds Differently}
\label{sec:composition}

The results so far describe the market as a whole. This section asks who the additional entrants are: independent, single-location operators, or established chains adding another outlet. Figure~\ref{fig:fact3dynamic} and Table~\ref{tab:fact3} report the dynamic version of equation~\eqref{eq:composition}, with a separate nonlegal-vs.\ always-legal-border coefficient $\hat\beta_{\tau,q}$ for each organizational type $\tau\in\{I,C\}$ and quarter $q$, together with a direct test of the independent-minus-chain gap in each quarter.

\textbf{Chains respond to captive cross-border demand significantly faster than independents in the first year after legalization; independents fully catch up within about a year, and from then on the two organizational forms are statistically indistinguishable.} At legalization (2022Q2), the chain coefficient is already positive and significant ($\hat\beta_C=+0.219$, $p<0.01$) while the independent coefficient is small and insignificant ($\hat\beta_I=+0.055$, $p=0.49$); the independent-minus-chain gap is significantly negative in each of the first four quarters (2022Q2--2023Q1, ranging from $-0.17$ to $-0.08$ log points, $p<0.05$ in three of four). By 2023Q2 independents have caught up, and for the remaining thirteen of seventeen quarters the gap is statistically indistinguishable from zero ($p$ ranging from $0.22$ to $0.98$), even as both coefficients continue climbing to a common plateau around $+0.32$ to $+0.39$.

This dynamic pattern explains why static, single-coefficient specifications give an unstable-looking picture of composition: a static regression can only report one number per outcome, which is mechanically close to the simple time-average of the quarter-specific effects (verified directly: averaging the seventeen $\hat\beta_{I,q}$ and $\hat\beta_{C,q}$ estimates reproduces the pooled full-panel static coefficients of $+0.304$ and $+0.339$ to three decimal places). That average cannot distinguish a small, persistent difference between organizational forms from a large, transient one that fully closes -- and here it is the latter: the composition gap the paper is built to explain is a first-year timing advantage for chains, not a lasting structural asymmetry in which type responds more to captive demand.

\begin{figure}[htbp]
\centering
\includegraphics[width=\textwidth]{fig_fact3_dynamic.png}
\caption{Fact 3, dynamic specification: quarter-specific nonlegal- vs.\ always-legal-border coefficients $\hat\beta_{\tau,q}$ for the number of independent and chain dispensaries. Sample: full New Mexico ZIP-code panel, 2022Q2--2026Q2 (368 ZIP codes, zero-inclusive). Inverse-hyperbolic-sine units; all quarters include $\log(\text{population})$ and calendar-quarter fixed effects; 95\% confidence intervals from county-clustered standard errors.}
\label{fig:fact3dynamic}
\end{figure}

\begin{table}[htbp]
\centering\scriptsize
\begin{tabular}{@{}lcccccc@{}}
\toprule
Quarter & \multicolumn{2}{c}{N independent brands (1)} & \multicolumn{2}{c}{N chain brands (2)} & \multicolumn{2}{c}{Indep.\ $-$ chain (3)}\\
\cmidrule(lr){2-3}\cmidrule(lr){4-5}\cmidrule(lr){6-7}
& Coef. & (SE) & Coef. & (SE) & Coef. & (SE) \\
\midrule
2022Q2 & +0.055 & (0.078) & +0.219*** & (0.079) & -0.173*** & (0.036) \\
2022Q3 & +0.127* & (0.066) & +0.243*** & (0.078) & -0.124*** & (0.035) \\
2022Q4 & +0.204*** & (0.071) & +0.280*** & (0.075) & -0.084* & (0.043) \\
2023Q1 & +0.240*** & (0.067) & +0.321*** & (0.079) & -0.089** & (0.038) \\
2023Q2 & +0.294*** & (0.061) & +0.342*** & (0.077) & -0.056 & (0.046) \\
2023Q3 & +0.331*** & (0.058) & +0.343*** & (0.085) & -0.020 & (0.043) \\
2023Q4 & +0.357*** & (0.061) & +0.358*** & (0.082) & -0.009 & (0.037) \\
2024Q1 & +0.371*** & (0.062) & +0.367*** & (0.080) & -0.004 & (0.029) \\
2024Q2 & +0.365*** & (0.064) & +0.372*** & (0.087) & -0.015 & (0.042) \\
2024Q3 & +0.380*** & (0.066) & +0.373*** & (0.085) & -0.001 & (0.040) \\
2024Q4 & +0.364*** & (0.059) & +0.387*** & (0.083) & -0.030 & (0.044) \\
2025Q1 & +0.371*** & (0.066) & +0.375*** & (0.084) & -0.012 & (0.048) \\
2025Q2 & +0.345*** & (0.063) & +0.380*** & (0.086) & -0.043 & (0.051) \\
2025Q3 & +0.372*** & (0.062) & +0.367*** & (0.085) & -0.003 & (0.052) \\
2025Q4 & +0.335*** & (0.060) & +0.359*** & (0.077) & -0.032 & (0.038) \\
2026Q1 & +0.334*** & (0.059) & +0.343*** & (0.076) & -0.017 & (0.042) \\
2026Q2 & +0.316*** & (0.067) & +0.337*** & (0.075) & -0.028 & (0.038) \\
\bottomrule
\end{tabular}

\caption{Fact 3, dynamic specification: quarter-specific coefficients for independent brands, chain brands, and their difference (independent minus chain), nonlegal- vs.\ always-legal-border ZIP codes. Same sample and controls as Figure~\ref{fig:fact3dynamic}; standard errors (in parentheses) clustered by county. $^{*}p<.10$ $^{**}p<.05$ $^{***}p<.01$.}
\label{tab:fact3}
\end{table}

\begin{table}[htbp]
\centering\footnotesize
\begin{tabular}{@{}lcccc@{}}
\toprule
 & \multicolumn{2}{c}{Main: nonlegal- vs.\ legal-border} & \multicolumn{2}{c}{Pooled: legal+interior pooled} \\
\cmidrule(lr){2-3}\cmidrule(lr){4-5}
Outcome (log) & Coef. & (SE) & Coef. & (SE) \\
\midrule
log(Total revenue) & +0.637** & (0.288) & +0.698*** & (0.175) \\
log(N stores) & +0.451** & (0.177) & +0.407*** & (0.100) \\
log(Revenue per store) & +0.186 & (0.175) & +0.291** & (0.120) \\
\bottomrule
\end{tabular}

\caption{Secondary conditional-on-activity specification: total revenue, dispensary count, and revenue per store (logs, market-age control), restricted to the 125 ZIP codes with at least one recorded dispensary. $^{*}p<.10$ $^{**}p<.05$ $^{***}p<.01$.}
\label{tab:fact1fact2conditional}
\end{table}

\begin{table}[htbp]
\centering\footnotesize
\begin{tabular}{@{}lcccc@{}}
\toprule
 & \multicolumn{2}{c}{Main: nonlegal- vs.\ legal-border} & \multicolumn{2}{c}{Pooled: legal+interior pooled} \\
\cmidrule(lr){2-3}\cmidrule(lr){4-5}
Outcome & Coef. & (SE) & Coef. & (SE) \\
\midrule
N independent brands & +1.103*** & (0.375) & +1.344*** & (0.416) \\
N chain brands & +1.038** & (0.475) & +0.638 & (0.434) \\
\bottomrule
\end{tabular}

\caption{Secondary conditional-on-activity specification: number of independent and chain dispensaries (levels, market-age control), same 125-ZIP-code sample as Table~\ref{tab:fact1fact2conditional}. $^{*}p<.10$ $^{**}p<.05$ $^{***}p<.01$.}
\label{tab:fact3conditional}
\end{table}

\begin{figure}[htbp]
\centering
\includegraphics[width=0.75\linewidth]{fig_fact3_composition.png}
\caption{Chain and independent store counts against a continuous captive-demand index $D(B)$ and its always-legal-border placebo $D(B_{\text{legal}})$, ZIP-quarter panel, pooled-control sample (95\% CI). A companion specification to Table~\ref{tab:fact3}: same sample and comparison groups, with the discrete $\text{Nonlegal}_z$ indicator replaced by the continuous demand indices, from this project's earlier ZIP-level structural work, predating the dynamic decomposition reported below. Read the apparent independents-only pattern shown here alongside Figure~\ref{fig:fact3dynamic}: pooling every quarter into one coefficient, as this exercise does, is exactly the kind of specification that understates the chain response, since chains reach most of their eventual response in the first quarter and contribute comparatively little identifying variation to a pooled estimate averaged over four subsequent years.}
\label{fig:fact3composition}
\end{figure}

\textbf{This pattern is not uniform across every specification, and the paper reports that plainly rather than presenting only the cleanest version.} The town-level aggregation (Appendix Table~\ref{tab:townlevel}), a full zero-inclusive panel like the main ZIP-code specification, finds independent entrants significant in both the main and pooled-control samples ($p<0.05$ or better) while chain entrants are not significant at the town level in either; the conditional-on-activity secondary specification (Appendix Table~\ref{tab:fact3conditional}) shows the same pattern. Figure~\ref{fig:fact3dynamic} points to a likely reason: both of these specifications compress every quarter into a single static coefficient, and a static coefficient is only as informative as the process it averages over. Chains reach most of their eventual response in the first year and change comparatively little afterward, so they contribute little of the variation a static, all-quarters estimate relies on; independents are still visibly rising throughout the same window and so contribute more of it. A coarser level of aggregation or a sample restricted to already-active markets can tilt that averaging further away from the early window in which chains actually move, making the chain coefficient look smaller and less precise than the dynamic pattern above shows it to be. This -- a real difference in \emph{when} each type responds, read by a static specification as a difference in \emph{whether} it responds at all -- is the empirical fact the entry game of Section~\ref{sec:model} is built to explain, once the sample is restricted to the settled, post-transition period the model targets.

The difference in entrant composition is economically important independent of its statistical precision in any one specification: whether a market's growth in response to demand comes from independent or chain operators has different implications for local ownership, price competition, and market power, and a market-size argument alone -- more demand, more of any kind of firm -- does not by itself explain why the two organizational forms would respond at different rates in the first place. Section~\ref{sec:model} develops an entry-threshold model with two ex ante heterogeneous pools of potential entrants and a specific, testable candidate mechanism for the composition asymmetry, and Section~\ref{sec:counterfactuals} reports the direct test of that mechanism, including where it fails to confirm the story.

\FloatBarrier

\section{A Strategic Entry Game for ZIP-Code Markets}
\label{sec:model}

Section~\ref{sec:composition} establishes that independent and chain potential entrants do not respond uniformly to captive cross-border demand, in two distinct ways: they respond on different timelines immediately after legalization, and, once the market has settled, an entry-threshold game finds their steady-state profitability responses differ as well. This section builds a discrete entry-threshold game for the second fact, in the tradition of \citet{BresnahanReiss1991}, \citet{Mazzeo2002}, and \citet{Seim2006}, in which two ex ante heterogeneous pools of potential entrants -- independent and chain -- each decide whether to enter a ZIP-code market, and derives a primitive under which the two pools' profitability responds differently to the same cross-border-demand shock. Section~\ref{sec:diffusion} returns to the first fact with a lightweight dynamic model of its own, designed to characterize the transition rather than to replace the steady-state game. Unlike Seim's or Mazzeo's applications, in which potential entrants choose a location or a product type but are otherwise symmetric, the two pools here are asymmetric by construction: only the chain pool's entrants can already have a store nearby. This is the same margin that \citet{Nishida2008} and \citet{Holmes2011} model directly for a single chain's own network-expansion decision; Section~\ref{sec:cannibalization} adapts it to a market-level entry game with two competing organizational types.

\subsection{Economic environment}
\label{sec:environment}

The unit of analysis is a ZIP-code market $z$ in quarter $t$. Two ex ante heterogeneous pools of potential entrants exist in every market-period: pool $I$ (independent, first-location entrants) and pool $C$ (established, multi-store operators). Within each pool, potential entrants are ranked $k=1,2,\dots$, a standard Bresnahan--Reiss/Mazzeo entry-threshold device rather than a literal arrival order. The $k$-th independent potential entrant's profit if it enters, given $N^C_{zt}$ active chain rivals and $k-1$ active independent rivals ahead of it in the ranking, is
\begin{equation}
\pi^I_{zt}(k) \;=\; \kappa_I \;+\; \beta_I \log L_{zt} \;+\; \delta_I \log(1+B_{zt}) \;-\; \gamma_{II}(k-1) \;-\; \gamma_{IC} N^C_{zt} \;+\; \varepsilon^I_{k,zt}. \label{eq:piI}
\end{equation}
Symmetrically, the $k$-th chain potential entrant's profit is
\[
\pi^C_{zt}(k) \;=\; \kappa_C \;+\; \beta_C \log L_{zt} \;+\; \delta_C \log(1+B_{zt}) \;-\; \gamma_{CC}(k-1) \;-\; \gamma_{CI} N^I_{zt} \;+\; \varepsilon^C_{k,zt}.
\]
$L_{zt}$ is ZIP-code $z$'s resident population, the same local-demand primitive as $\text{Population}_z$ in equation~\eqref{baseline}; $B_{zt}$ is the captive-cross-border-demand index, the population living within 75 miles across a nonlegal-state line, zero for always-legal-border and interior ZIP codes. $\kappa_\tau$ ($\tau\in\{I,C\}$) is a type-specific net intercept identifying profitability net of entry cost, not entry cost separately, absent a cost shifter that varies within type. $\delta_\tau$ is the parameter of central interest -- the type-specific cross-border-demand response, the model's counterpart to $\hat\beta_I$ and $\hat\beta_C$ in Table~\ref{tab:fact3} -- and $\gamma_{\tau\tau'}$ are competitive effects.

\paragraph{Competitive-effect structure.} A fully saturated specification would let all four of $\gamma_{II}, \gamma_{CC}, \gamma_{IC}, \gamma_{CI}$ differ. Section~\ref{sec:identification} shows the New Mexico town-market panel does not support this: conditional on $\log L_{zt}$, $\log(1+B_{zt})$, and quarter effects, $N^I_{zt}$ and $N^C_{zt}$ retain a substantial residual correlation, because both are driven by the same local-demand primitives. The preferred specification therefore imposes a parsimonious same/other-type restriction, $\gamma_{II}=\gamma_{CC}\equiv\gamma_{\text{same}}$, $\gamma_{IC}=\gamma_{CI}\equiv\gamma_{\text{other}}$: a potential entrant's profit falls by $\gamma_{\text{same}}$ per same-type rival and $\gamma_{\text{other}}$ per other-type rival. The unrestricted four-parameter model is estimated in parallel as an identification diagnostic only, in keeping with \citet{BresnahanReiss1991}'s caution against over-parameterizing entry-threshold models beyond what the cross-section of markets supports.

\subsection{Own-network internalization: why $\delta_I$ may exceed $\delta_C$}
\label{sec:cannibalization}

The specification above lets $\delta_I$ and $\delta_C$ differ freely, and Table~\ref{tab:fact3} and Section~\ref{sec:modelfit} show the data want them to. Left there, that inequality is only a fitted fact -- consistent with many stories, including one in which chains simply have a flatter demand response for reasons unrelated to their own network. This subsection derives $\delta_I \neq \delta_C$ from a primitive that does not, by construction, presuppose the answer, adapting the network-cannibalization logic of \citet{Nishida2008} and \citet{Holmes2011} to a market with two competing organizational types rather than one chain's own expansion problem.

The two organizational types face the same primitive cross-border-demand expansion: a more exposed location draws more gross revenue, whatever type opens it. They differ in how much of that gross revenue is incremental to the firm. A first-location independent entrant has no other New Mexico outlet, so every dollar the new store draws is a dollar of new firm-wide revenue. An established multi-store operator opening an additional outlet may instead partly relocate demand from its own nearby stores: a customer who would have driven to the operator's existing location now visits the new one instead, so part of the new store's revenue is a transfer within the firm rather than a firm-wide increase. This is captured with a predetermined network-overlap state, $H_{jzt}\ge 0$, observed and computed only from operator $j$'s storefronts already active in the quarter before $t$, never from stores that open later.

Let $D(B_z)\equiv\log(1+B_z)$ and let $\delta>0$ denote the common, primitive gross-profitability effect of cross-border demand on a marginal storefront, before any own-network offset. The $k$-th independent potential entrant's profit is unchanged from equation~\eqref{eq:piI}, since its own-network overlap is zero by construction. An established operator $j$'s marginal outlet in ZIP code $z$ at $t$ has profit
\begin{equation}
\pi^C_{jzt} \;=\; \kappa_C + \beta_C\log L_{zt} + \big[\delta - \rho H_{jzt}\big]D(B_z) - \text{competition} + \varepsilon^C_{jzt}, \label{eq:piC_micro}
\end{equation}
where $\rho\ge 0$ is the strength of internalized own-network cannibalization. Differentiating,
\[
\frac{\partial \pi^I_{zt}}{\partial D(B_z)} = \delta, \qquad \frac{\partial \pi^C_{jzt}}{\partial D(B_z)} = \delta - \rho H_{jzt}, \qquad \Longrightarrow \qquad \frac{\partial \pi^I_{zt}}{\partial D(B_z)} - \frac{\partial \pi^C_{jzt}}{\partial D(B_z)} = \rho H_{jzt} \;\ge\; 0,
\]
strictly positive whenever $\rho>0$ and $H_{jzt}>0$.

\begin{proposition}[Organizational heterogeneity from own-network internalization]
\label{prop:cannibalization}
Suppose an established multi-store operator internalizes the effect of a new outlet on its own pre-existing network ($\rho>0$), and the candidate market overlaps that network ($H_{jzt}>0$). Then the marginal profitability response to cross-border demand of a first-location independent entrant is at least as large as that of the established operator's marginal outlet, and strictly larger whenever $H_{jzt}>0$:
\[
\frac{\partial \pi^I_{zt}}{\partial D(B_z)} \;\ge\; \frac{\partial \pi^C_{jzt}}{\partial D(B_z)}, \qquad \text{with equality iff } \rho H_{jzt}=0.
\]
\end{proposition}
\begin{proof}
Immediate from the two derivatives above: their difference is $\rho H_{jzt}$, non-negative by $\rho\ge0, H_{jzt}\ge0$, and strictly positive iff both are strictly positive.
\end{proof}

This is a statement about the marginal network profit of an additional outlet, not a claim that independent firms are more productive, better-run, or more profitable than chains on average -- the model is silent on, and does not need, any such claim. Mapped to the reduced-form parameters of Section~\ref{sec:environment}: $\delta_I=\delta$ is the independent pool's response, since no operator has a positive $H$ before its first store, while $\delta_C$ is what a market-level model that does not carry an operator index recovers as an average of $\delta-\rho H_{jzt}$ across established operators actually expanding. Section~\ref{sec:counterfactuals} constructs $H_{jzt}$ from this project's own storefront coordinates and tests Proposition~\ref{prop:cannibalization}'s prediction directly at the operator level, rather than treating the aggregate $\hat\delta_I>\hat\delta_C$ finding as proof of a mechanism that could equally be consistent with other, unmodeled sources of type-specific demand response.

\subsection{Private information and equilibrium}
\label{sec:equilibrium}

The private shocks $\varepsilon^I_{k,zt}, \varepsilon^C_{k,zt}$ are i.i.d.\ standard logistic across potential entrants, types, markets, and periods, a scale normalization rather than an estimated parameter, standard in this literature \citep{Seim2006,TrainDCM}. A potential entrant enters iff its realized profit is positive, so $\Pr(\text{entry}\mid N^{-\tau}_{zt}) = \Lambda(\pi^\tau_{zt}(k))$ with $\Lambda$ the logistic CDF. With a pool of ex ante symmetric potential entrants per type, capped at $\bar K=80$ slots, comfortably above the largest observed market entrant count, the expected number of type-$\tau$ entrants given a belief about the rival type's count is
\[
R_I(N^C_{zt}) = \sum_{k=1}^{\bar K}\Lambda(\pi^I_{zt}(k)), \qquad R_C(N^I_{zt}) = \sum_{k=1}^{\bar K}\Lambda(\pi^C_{zt}(k)).
\]
A \textbf{Bayesian--Nash equilibrium} of the market-period $(z,t)$ game is a pair $(N^{I*}_{zt}, N^{C*}_{zt})$ such that each type's expected entrant count is a best response to the other's: $N^{I*}_{zt}=R_I(N^{C*}_{zt})$, $N^{C*}_{zt}=R_C(N^{I*}_{zt})$. Because $\gamma_{\text{same}},\gamma_{\text{other}}\ge0$ (Section~\ref{sec:modelfit}), the game is one of strategic substitutes. The fixed point is solved numerically by a damped-fixed-point/Newton hybrid (Appendix~\ref{app:computation}); Section~\ref{sec:modelfit} reports whether multiple starting points converge to the same fixed point at the estimated parameters, separately for the ZIP-code and town-level estimations, since this diagnostic differs between the two.

\begin{proposition}[Demand expands entry]
\label{prop:1}
Holding $L_{zt}$ fixed, $N^{I*}_{zt}+N^{C*}_{zt}$ is weakly increasing in $B_{zt}$ whenever $\delta_I,\delta_C\ge0$: both best-response curves shift out in $B_{zt}$, and a system of decreasing, upward-shifting best responses has a weakly larger fixed point. This is the model's counterpart to P1.
\end{proposition}

\begin{proposition}[Heterogeneous response implies composition shifts with demand]
\label{prop:2}
Under the microfoundation of Section~\ref{sec:cannibalization} with $\rho>0$ and $H_{jzt}>0$ for prospective chain expansions, Proposition~\ref{prop:cannibalization} gives $\partial\pi^I/\partial D(B_z)\ge\partial\pi^C_j/\partial D(B_z)$ at the level of a single potential entrant's payoff, so each independent potential entrant's entry probability rises at least as fast in $D(B_z)$ as each chain potential entrant's does. Whether this translates into the \emph{equilibrium} composition $N^{I*}_{zt}/(N^{I*}_{zt}+N^{C*}_{zt})$ rising in $B_{zt}$ additionally requires that the equilibrium feedback through $\gamma_{\text{same}},\gamma_{\text{other}}$ not overturn the direct effect; a sufficient condition is $\gamma_{\text{same}}=\gamma_{\text{other}}$, the restriction the data support (Section~\ref{sec:identification}). This is the model's counterpart to P3.
\end{proposition}

These are theoretical predictions, verified against the estimated model in Sections~\ref{sec:modelfit}--\ref{sec:counterfactuals}, not themselves empirical findings.

\subsection{A dynamic extension: readiness, not speed}
\label{sec:diffusion}

The entry-threshold game above is a model of the steady state $(N^{I*}_{zt}, N^{C*}_{zt})$ that a market settles into once entrants have had time to act on a demand shock; it has nothing to say about how a market gets there. Figure~\ref{fig:fact3dynamic} shows that the path matters: chain and independent entry counts start far apart at legalization and converge over roughly a year. This subsection asks what a market's transition path looks like, using the same seventeen quarter-specific coefficients $\hat\beta_{\tau,q}$ from Table~\ref{tab:fact3}, and finds an answer sharper than ``chains move faster.''

\paragraph{A common-speed, different-readiness model.} Suppose each organizational type's eventual, steady-state response to captive demand, $\beta_\tau^\infty$, is the same object the entry-threshold game characterizes -- common across types once the market has settled, consistent with the entry-threshold game finding independent and chain profitability responses that are close relative to the precision of a dynamic, quarter-by-quarter estimate. What differs by type is how much of that eventual response is already realized at the moment the market opens, and how quickly the remainder is realized afterward. Let $\phi_\tau\in[0,1]$ be the share of type $\tau$'s steady-state response already present at legalization ($t=0$), and let the remaining share close at a rate $\lambda_\tau$ common to both types if a single institutional bottleneck -- for instance, the pace at which the state's own licensing process clears applications -- governs how fast \emph{any} applicant, independent or chain, can convert a decision to enter into an open store. This gives
\begin{equation}
\label{eq:diffusion}
\beta_{\tau,q} = \beta^\infty - \big(\beta^\infty - \beta^\infty\phi_\tau\big)e^{-\lambda_\tau q}, \qquad \tau\in\{I,C\},
\end{equation}
with $q$ measured in quarters since legalization. $\phi_\tau=1$ says type $\tau$ needs no transition at all; $\phi_\tau=0$ says it starts from nothing and closes the entire gap at rate $\lambda_\tau$, the pure diffusion case.

\paragraph{Estimates.} Fitting equation~\eqref{eq:diffusion} to the seventeen quarterly coefficients underlying Table~\ref{tab:fact3} and Figure~\ref{fig:fact3dynamic} by nonlinear least squares, weighted by the inverse variance of each $\hat\beta_{\tau,q}$, gives a shared long-run response $\hat\beta^\infty=+0.363$ (SE $0.018$) and a shared adjustment speed $\hat\lambda=0.408$ (SE $0.148$), implying a half-life of $1.7$ quarters for whichever share of a type's response has not yet materialized. Imposing $\lambda_I=\lambda_C$ costs almost nothing in fit (the weighted sum of squared residuals rises from $2.514$ to $2.517$ on $29$ degrees of freedom, far short of any conventional rejection threshold), so the data do not merely fail to reject a common speed -- they actively prefer one. What differs sharply by type is readiness: $\hat\phi_C=0.57$, so chains arrive at legalization already $57$ percent of the way to their eventual response, against $\hat\phi_I=0.09$ for independents. Figure~\ref{fig:diffusionfit} plots the fitted curves against the underlying coefficients, and Table~\ref{tab:diffusionfit} reports the estimates.

\begin{figure}[htbp]
\centering
\includegraphics[width=\textwidth]{fig_fact3_diffusion_fit.png}
\caption{Common-speed, different-readiness diffusion model (equation~\eqref{eq:diffusion}) fitted to the dynamic composition coefficients of Figure~\ref{fig:fact3dynamic}. Points are the quarter-specific $\hat\beta_{\tau,q}$ estimates with 95\% confidence intervals; lines are the fitted curves under a shared adjustment speed $\hat\lambda$ and type-specific readiness shares $\hat\phi_I,\hat\phi_C$.}
\label{fig:diffusionfit}
\end{figure}

\begin{table}[htbp]
\centering\footnotesize
\begin{tabular}{@{}lccc@{}}
\toprule
Parameter & Independent & Chain & Shared \\
\midrule
Long-run gap $\hat\beta_\infty$ & \multicolumn{2}{c}{+0.363 (0.018)} & \\
Initial gap $\hat\beta_{\tau,0}$ & +0.032 (0.067) & +0.207 (0.060) & \\
Share of long-run gap reached at $t=0$, $\hat\phi_\tau$ & 0.087 & 0.570 & \\
Adjustment speed $\hat\lambda$ (shared) & \multicolumn{2}{c}{+0.408 (0.148)} & \\
Half-life (quarters) & \multicolumn{2}{c}{1.70} & \\
\bottomrule
\end{tabular}

\caption{Fitted parameters of the common-speed, different-readiness diffusion model, equation~\eqref{eq:diffusion}. Standard errors in parentheses, from weighted nonlinear least squares (weights inverse to each quarterly coefficient's estimated variance).}
\label{tab:diffusionfit}
\end{table}

\paragraph{Reading the result.} A common adjustment speed is what a shared regulatory bottleneck predicts: the Cannabis Control Division processes every application on the same clock, whoever files it. A large gap in $\phi_\tau$ is what a difference in pre-existing readiness predicts: an operator that has already raised capital, scouted real estate, and filed licenses in other states can have a large share of its eventual footprint already in the application pipeline before a captive-demand opportunity is even public knowledge, while a first-time independent has to complete each of those steps from scratch. This is a different candidate explanation from the own-network-internalization mechanism of Section~\ref{sec:cannibalization}: cannibalization is a reason $\delta_C$ could be permanently smaller than $\delta_I$, and it says nothing about why a gap would appear in quarter one and be gone within a year. A readiness gap that is worked off at a common, roughly institutional speed does exactly that, and it does so without requiring a new profitability parameter -- it is a statement about the timing of execution, not about the size of the opportunity either type ultimately captures. Read together with the null result on own-network internalization in Section~\ref{sec:mechreg}, the paper now has two independent pieces of evidence pointing away from cannibalization as the operative channel for \emph{why} composition responds to captive demand at all, even as the entry-threshold game's finding that steady-state profitability responses differ once the market has settled stands on its own.

This exercise is deliberately not a re-estimation of the Bayesian--Nash entry game in event time: doing so would require a dynamic equilibrium concept, a belief-formation process during the transition, and considerably more data per market-period than seventeen aggregate quarters provide. Equation~\eqref{eq:diffusion} is a reduced-form description of the transition path chosen for its transparency and its tight fit, not a claim to have modeled the strategic interaction between potential entrants while the market is still adjusting; Section~\ref{sec:limitations} states this scope restriction directly.

\FloatBarrier

\section{Identification and Estimation}
\label{sec:identification}

\subsection{Sample and unit of analysis}

Consistent with the rest of the paper, the model is estimated primarily at the ZIP-code level: 97 ZIP codes that record at least one dispensary at some point in the panel, restricted to each ZIP code's own \emph{converged} quarters -- market age at least eight quarters since that ZIP code's first active dispensary -- so that the observed $(N^I_{zt}, N^C_{zt})$ is plausibly close to its stationary Bayesian--Nash equilibrium rather than still transitioning from the market's opening. This threshold reuses the finding, established elsewhere in this project's town-level work, that chain/independent composition converges within roughly six to eight quarters of a market's first dispensary, a window consistent with the estimated $1.7$-quarter adjustment half-life in Section~\ref{sec:diffusion}: after five to six half-lives, the transition is effectively complete. Albuquerque's constituent ZIP codes are excluded from the ZIP-code estimation, for the same reason a single very large, atypical market is excluded from the town-level estimation below: including it materially worsens both model fit and the stability of the equilibrium solver, documented directly in Section~\ref{sec:modelfit} rather than assumed. The resulting ZIP-code estimation sample has 81 ZIP codes and several hundred ZIP-quarters, reported exactly in Table~\ref{tab:zipstructural}.

The \textbf{town-level} estimation, reported throughout as a robustness check per this paper's stated design, aggregates the same underlying storefront data to 78 town-markets (again excluding Albuquerque), using an analogous converged-quarter restriction. This is the unit used in an earlier version of this project's structural work and is retained here specifically to check whether the ZIP-code result depends on treating a ZIP code, rather than a town, as the relevant competitive market.

\subsection{Which variation identifies each parameter}

$\beta_\tau$ is identified from cross-sectional variation in $\log L_{zt}$ conditional on $B_{zt}$ and quarter effects. $\delta_\tau$, the parameter of central interest, is identified from variation in $B_{zt}$ conditional on $\log L_{zt}$ and quarter effects: markets close to a nonlegal-state line have $B_{zt}>0$, markets near an always-legal-state line or in the interior have $B_{zt}\approx0$, and the model attributes the resulting difference in $N^{I*}_{zt}$ versus $N^{C*}_{zt}$ separately to $\delta_I$ and $\delta_C$. $\kappa_\tau$ is identified from the level of entry once $L_{zt}$ and $B_{zt}$ are accounted for, and is a composite of profitability and entry cost, not separately identified into the two without a cost shifter that varies within type -- a limitation stated plainly in Section~\ref{sec:limitations}. $\gamma_{\text{same}}, \gamma_{\text{other}}$ are identified from how each type's entrant count responds to the other type's realized count, conditional on $L_{zt}$ and $B_{zt}$; the same/other restriction (Section~\ref{sec:environment}) is imposed because, at both the ZIP-code and town-level sample sizes, $N^I_{zt}$ and $N^C_{zt}$ retain enough residual correlation after conditioning on $L_{zt}$, $B_{zt}$, and quarter effects that a fully saturated four-parameter competitive-effect structure is not well supported; the unrestricted estimate is reported as a diagnostic only.

\subsection{Estimator}

The model is estimated by simulated minimum distance: for a candidate parameter vector $\theta$, the Bayesian--Nash fixed point $(N^{I*}_{zt}(\theta), N^{C*}_{zt}(\theta))$ is solved for every market-period (Appendix~\ref{app:computation}), and $\theta$ is chosen to minimize the weighted sum of squared deviations from the observed $(N^I_{zt}, N^C_{zt})$, with weights approximating the variance of a Poisson count. Multiple optimizer starting points are used and their spread reported (Section~\ref{sec:modelfit}) as a check against local optima; standard errors are computed by a market-clustered nonparametric bootstrap, clustering by ZIP code in the ZIP-code estimation and by town in the town-level estimation.

\section{Model Fit and Validation}
\label{sec:modelfit}

Re-estimating the identical restricted BNE model at ZIP resolution (81 ZIPs with a dispensary in the converged sample -- per-ZIP market age $\geq 8$ quarters, 625 ZIP-quarters -- versus 78 towns at town resolution) gives $\hat\delta_I=+0.047$ against $\hat\delta_C=+0.002$ ($\hat\delta_I-\hat\delta_C=+0.046$, boot.\ SE 0.016, positive in 100\% of 10 ZIP-clustered bootstrap replications), the SAME SIGN as the town-level result ($\hat\delta_I=0.052$, $\hat\delta_C=0.009$). The identification concern raised before attempting this is confirmed, not merely feared: residual corr($N_I$,$N_C$) after partialling out $\log(\text{pop})$, $\log(1+B)$, and quarter fixed effects is 0.47 at ZIP resolution, versus 0.67 at town resolution -- comparable to the town-level figure, so this specific concern does not bind harder at ZIP resolution than it already did at town resolution. In-sample fit across ZIP means: corr(actual,model) = 0.80 for total stores, 0.78 for independent, 0.68 for chain (town-level benchmark: 0.93/0.93/0.83).

\begin{table}[htbp]
\centering\footnotesize
\begin{tabular}{@{}lrrr@{}}
\toprule
Parameter & Estimate & Boot.\ SE & 95\% CI \\
\midrule
kappa\_I & -0.9506 & (0.1764) & [-1.043, -0.510] \\
kappa\_C & -2.3644 & (0.2283) & [-2.964, -2.363] \\
beta\_I(logpop) & +0.9293 & (0.1477) & [+0.636, +1.060] \\
beta\_C(logpop) & +0.0000 & (0.0224) & [+0.000, +0.063] \\
delta\_I(logB75) & +0.0472 & (0.0145) & [+0.023, +0.071] \\
delta\_C(logB75) & +0.0017 & (0.0040) & [+0.000, +0.012] \\
gamma\_same & +0.3135 & (0.0396) & [+0.226, +0.326] \\
gamma\_other & -0.3720 & (0.0254) & [-0.393, -0.314] \\
\bottomrule
\end{tabular}

\caption{Estimated restricted entry game, ZIP-code level (primary estimation, Albuquerque excluded, converged quarters). $\hat\delta_I$ and $\hat\delta_C$ are the type-specific cross-border-demand responses; $\hat\delta_I-\hat\delta_C$ is the model's counterpart to the composition asymmetry documented in Table~\ref{tab:fact3}.}
\label{tab:zipstructural}
\end{table}

\begin{figure}[htbp]
\centering
\includegraphics[width=0.9\linewidth]{fig_zipstructural_fit.png}
\caption{ZIP-code entry-game fit: actual vs.\ model-predicted store counts, converged-quarter ZIP-code means, Albuquerque excluded.}
\label{fig:zipstructuralfit}
\end{figure}

The town-level estimation, reported as a robustness check per this paper's design, gives $\hat\delta_I=0.052$ against $\hat\delta_C=0.009$ ($\hat\delta_I-\hat\delta_C=+0.042$, positive in 100\% of bootstrap replications), with a uniquely converged Bayesian--Nash fixed point from every tested starting point, and model-versus-data rank correlations of 0.93 (total stores), 0.93 (independent), and 0.83 (chain) across town means. The ZIP-code point estimate agrees on sign and is, if anything, similar or slightly larger in magnitude ($\hat\delta_I-\hat\delta_C=+0.046$), but comes from a less well-behaved equilibrium: the fixed point is not unique from every starting point tested, and fit correlations are lower (Figure~\ref{fig:zipstructuralfit}), consistent with a ZIP code being a noisier, more thinly populated approximation to a competitive market than a town is. Both estimations therefore agree on the qualitative finding this paper's structural work is built to explain -- $\hat\delta_I>\hat\delta_C$ -- and the town-level estimation is used for the counterfactual exercise below specifically because it is the one with a reliably unique equilibrium to re-solve.

\FloatBarrier

\section{Counterfactuals and the Own-Network Mechanism}
\label{sec:counterfactuals}

\subsection{Legalization counterfactual}
\label{sec:legalization_cf}

To gauge the economic magnitude of $\hat\delta_I>\hat\delta_C$, the estimated town-level model -- used here because its uniquely converged equilibrium (Section~\ref{sec:modelfit}) supports a stable counterfactual re-solve, a property the ZIP-code estimation does not share -- is used to ask what would happen to entry in nonlegal-border town-markets if the neighboring nonlegal states legalized recreational cannabis, so that $B_{mt}\to0$ for the 28 town-markets currently exposed to nonlegal-state demand, holding population, geography, and border status otherwise fixed. Re-solving the Bayesian--Nash equilibrium under this counterfactual gives $\Delta N_I = -518.0$ and $\Delta N_C = -112.4$ (95\% bootstrap confidence intervals excluding zero for both): 82.2 percent of the net cross-border-demand increment in the town-level model is independent entry. This is the town-level model's quantitative counterpart to the composition fact documented directly in Table~\ref{tab:fact3}: not only does independent entry respond more per unit of cross-border demand, it accounts for the large majority of the store count that a captive-demand shock of this size adds to or removes from the market.

\subsection{Testing the own-network mechanism directly}
\label{sec:mechreg}

Proposition~\ref{prop:cannibalization} makes a testable prediction: an established operator's cross-border-demand response should weaken as its own pre-existing network overlaps more with the candidate market. This section builds the network-overlap state $H_{jzt}$ from data already in this project and tests the prediction directly at the operator level, reporting the result whichever way it comes out.

Every New Mexico storefront carries a geocoded latitude and longitude (100 percent coverage). For every operator $j$, ZIP code $z$, and quarter $t$, using only storefronts of $j$ active in quarter $t-1$, $H^{\text{same}}_{jzt}$ counts $j$'s stores active in the same ZIP code at $t-1$ and $H^{\text{geo}}_{jzt}(\lambda)=\sum_{r} e^{-\lambda\,d(z,r)}$ sums an exponential distance decay over $j$'s existing stores, with $d(z,r)$ the great-circle distance to store $r$. Of 370 realized chain-expansion events, 62--82 percent occur where the operator already has a same-market or nearby store: own-network overlap is common among established operators, not rare, so the mechanism is not ruled out on the grounds of being empirically irrelevant.

The mechanism's own key prediction is tested with an operator-and-quarter fixed-effects PPML regression on a candidate-market risk-set panel (established operator $\times$ market $\times$ quarter, every market-period in which that operator did not yet have a store, several hundred thousand rows), estimated by a sparse two-way fixed-effects Poisson IRLS (Appendix~\ref{app:ppml}) since dense dummy encoding at this scale is not memory-feasible:
\[
E[\text{EntryC}_{jzt}\mid\cdot] = \exp\{\alpha_j + \alpha_t + \beta\log L_{zt} + \delta_B D(B_z) + \theta H_{jzt} + \varphi\,(H_{jzt}\times D(B_z))\}.
\]
The theoretically central coefficient is $\varphi$, predicted negative by Proposition~\ref{prop:cannibalization}: more own-network overlap should mean a weaker response to cross-border demand. \textbf{It is not.} In the preferred specification (established-operator sample, market and quarter fixed effects), $\hat\varphi_{H_{same}}=-0.014$ (SE 0.020, $p=0.49$) and $\hat\varphi_{H_{geo}}=+0.022$ (SE 0.018, $p=0.21$): both statistically indistinguishable from zero, and the $H_{geo}$ point estimate has the wrong sign. Including the largest metro market as a robustness check does not rescue the prediction: both interaction estimates become significantly \emph{positive}, the opposite of what own-network internalization predicts. \textbf{The direct test does not support $\varphi<0$ (equivalently $\rho>0$) in these data.}

\subsection{What the null result does, and does not, rule out}
\label{sec:mechanism-conclusion}

Three things are true simultaneously, and the paper states all three rather than the one that is most convenient. First, the aggregate pattern this section set out to explain, $\hat\delta_I>\hat\delta_C$, is robust across the ZIP-code and town-level estimations (Section~\ref{sec:modelfit}) and across every composition specification in Table~\ref{tab:fact3} and its robustness checks that finds independent entrants responding at all. Second, own-network overlap among established operators' expansions is common, so the own-network-internalization mechanism is not ruled out on the grounds of being empirically irrelevant. Third, and this is the qualification the paper does not soften elsewhere to compensate for: the direct operator-level test of the mechanism's own parameter, $\rho$, does not confirm it. \textbf{This paper therefore presents own-network internalization as a theoretically coherent, well-motivated candidate mechanism for $\hat\delta_I>\hat\delta_C$ -- not as an empirically identified one.}

A null result on $\rho$ does not mean the composition asymmetry is unexplained; it means this particular, testable channel is not the confirmed explanation. At least three other channels are consistent with everything reported in this paper and are not directly tested here, for reasons stated alongside each. \emph{Financing constraints}: an independent entrant financing a single, geographically concentrated demand opportunity may accept a narrower, more concentrated risk than a chain's capital-allocation process would tolerate for the same expected return, a hypothesis that would require internal financing or cost-of-capital data this project does not have access to. \emph{Local information advantages}: a first-time entrant choosing where to open at all may rely on local knowledge of a specific cross-border commuting pattern that a chain's centralized site-selection process, optimized across many candidate markets at once, does not weight as heavily; testing this would require data on how each type of operator actually searched for and evaluated candidate sites, which this project's licensee and storefront records do not contain. \emph{Differential risk tolerance toward a policy-contingent demand source}: captive cross-border demand is, by construction, contingent on a neighboring state's legalization status, and a chain with a diversified, primarily interior portfolio may discount a demand source that could disappear with a single neighboring-state policy change more heavily than a first-time entrant evaluating one location on its own terms; distinguishing this from the internalization story tested above would require either a natural experiment in a neighboring state's legalization timing, which the sample period does not contain, or direct evidence on operators' expectations, which is not observed. The paper does not attempt to adjudicate among these three, or the internalization mechanism itself, with the data at hand; it reports what was tested, what the test found, and states plainly that the composition fact in Table~\ref{tab:fact3} remains established while its underlying cause remains open.

\FloatBarrier

\section{Robustness}
\label{sec:falsification}

The main comparison, nonlegal- versus always-legal-border ZIP codes, is itself a robustness-motivated redesign relative to a simple border-versus-interior split: because both groups sit near a state line, it already rules out generic border proximity -- transportation access, land values, unmeasured local commercial activity -- as an explanation for Facts 1--3. This section reports the additional checks specific to the remaining concerns: whether the results depend on the particular distance threshold used to define exposure, whether the pooled-control, town-level, and conditional-on-activity samples change the picture, and whether cannabis-market activity near a nonlegal-state line reflects something broader about commercial activity in the same ZIP codes.

\textbf{Pooled control, conditional-on-activity, town level, and Albuquerque.} These four checks, introduced in Sections~\ref{sec:empirical}--\ref{sec:composition}, are reported together here for reference. Appendix Tables~\ref{tab:fact1fact2conditional} and~\ref{tab:fact3conditional} re-estimate the main specifications restricted to the 125 ZIP codes with recorded dispensary activity, in logs with a market-age control; Appendix Table~\ref{tab:townlevel} re-estimates the full zero-inclusive main specification at the town level; Appendix Table~\ref{tab:fact3abq} reports the composition result with Albuquerque included and excluded from the pooled-control sample. As discussed in Section~\ref{sec:composition}, the revenue and dispensary-density results (Facts 1--2) are stable across all of these; the composition result (Fact 3) is where these checks disagree most on the surface, though the disagreement has a single, well-defined source. Independent entrants respond significantly and positively in every specification. Chain entrants respond at least as strongly as independents in the full, zero-inclusive ZIP-code panel and less strongly, and not always significantly, once the sample is restricted to already-active markets or aggregated to the town level -- but the dynamic breakdown in Section~\ref{sec:composition} shows this is not four specifications disagreeing about whether chains respond; it is four specifications giving more or less weight to the first year after legalization, which is the only period in which chains and independents actually differ.

\textbf{Distance thresholds.} Border status is redefined using 50-mile and 100-mile thresholds in place of the baseline 75 miles, and dispensaries are assigned to 25-mile distance bands from the relevant state line. Figure~\ref{fig:eventstudy_50_100mi} and Figure~\ref{fig:nm_doseresponse} show that the main results are not an artifact of the specific 75-mile cutoff: point estimates are broadly similar across thresholds, with wider confidence intervals at 100 miles, as expected from the broader, less directly exposed sample it includes. Appendix Table~\ref{tab:nm_nonlegalpop_doseresponse} and Figure~\ref{fig:nm_nonlegalpop_doseresponse} show the premium also increases with the size of the nearby nonlegal-state population, consistent with a genuine demand-exposure gradient rather than a threshold artifact.

\textbf{Non-cannabis placebo sectors.} If nonlegal-border ZIP codes were simply more commercially active in general, that difference should also appear in unrelated retail sectors and should not track New Mexico's own cannabis-legalization date specifically. Figure~\ref{fig:noncannabis_eventstudy} interacts border status with each calendar quarter for alcohol and grocery retail activity, anchored to the legalization date. Neither sector shows a sharp change coinciding with cannabis legalization, evidence against a general commercial-activity confound.

Taken together, these checks do not rule out every possible source of confounding, nor is doing so the primary scope of the paper, but they narrow the plausible alternative explanations for Facts 1--3 considerably, and they are reported here, including where a check weakens rather than strengthens the main result, per the paper's general approach to reporting robustness honestly.

\FloatBarrier

\section{Discussion: Colorado and Constrained Entry}
\label{sec:discussion}

The New Mexico evidence shows how a relatively free-entry cannabis market absorbs captive cross-border demand: markets near a nonlegal-state line are larger than markets near an already-legal-state line, primarily through more dispensaries rather than higher revenue per store, and the additional entry is disproportionately independent. This is the pattern predicted by P1, P2, and P3: additional demand attracts more firms rather than raising the revenue of the stores already operating, and firms of different organizational forms respond at different rates \citep{BresnahanReiss1991}.

The paper next turns to P4 and asks whether the size of the density response depends on how entry is regulated. It compares two policy environments: New Mexico, where entry is largely unrestricted, and Colorado, where local governments can restrict or ban cannabis retailers. Both states border the same prohibition states, so both have counties exposed to a similar kind of cross-border demand. In the framework, New Mexico's density response to border exposure is the free-entry benchmark, and a smaller Colorado response would show that local entry barriers bind. The gap between the two responses is what P4 says can be read as the number of stores kept out and, applied to revenue per store, as the legal-market revenue and tax base forgone.

Appendix Table~\ref{tab:regime_interaction} reports the $\text{Border}_c \times \text{Restricted}_s$ interaction for all three outcomes, while Appendix Section~\ref{sec:crossstate} describes the county-level sample and specification. The interaction coefficient is negative for all three outcomes, so the border premium is smaller in Colorado than in New Mexico, though the estimates are imprecise. It is marginally significant for dispensary density ($\theta=-0.932$) and not statistically distinguishable from zero for revenue per capita ($\theta=-0.458$) or revenue per store. The direction still lines up with P4. In New Mexico, where entry is largely unrestricted, the premium is mostly a larger number of stores, an extensive-margin response. In Colorado, where entry is restricted, more of it shows up as higher revenue per store, the intensive-margin response the framework attributes to a binding cap.

To turn the density gap into dollars, the paper uses the dispensary-density estimate, the outcome for which the difference across entry regimes is statistically significant. The counterfactual asks how much more dispensary revenue Colorado's border counties would have generated in 2024 if their density response to border exposure had matched New Mexico's free-entry response rather than the restricted-entry one, holding revenue per store at its observed Colorado level. Applying the multiplier $e^{-\theta}\approx2.54$, based on $\theta=-0.932$, to each border county's own 2024 revenue and summing across counties, Colorado's border counties generated \$324.2 million in 2024, whereas the free-entry benchmark implies \$823 million, a difference of approximately \$499 million (Appendix Table~\ref{tab:counterfactual}). Because the interaction is only marginally significant and the bootstrap interval for this difference extends to zero, the dollar value should be read as an illustrative benchmark for the order of magnitude at stake rather than as a precise estimate of forgone revenue.

Because the entry regime is not randomly assigned and only two states are compared directly, the cross-state estimates are interpreted as suggestive evidence for P4 rather than as a causal estimate of the effect of licensing restrictions. Under this interpretation, the cross-state comparison complements the New Mexico results. The direction of the estimates is consistent with \citet{thomas2019license}, who shows that license restrictions can substantially constrain market expansion, and with evidence that entry restrictions in cannabis markets create local market power for incumbent retailers \citep{HollenbeckUetake2021}.

Taken together, the results suggest that whether cross-border demand expands a legal cannabis market or primarily raises revenue for existing stores, and whether that expansion is dominated by independent or chain entrants, depends on the entry environment and on organizational-form-specific profitability, both of which this paper documents but the second of which it cannot fully explain. For policymakers evaluating licensing and zoning rules, this distinction identifies an important margin: tighter entry restrictions may limit the number of dispensaries operating near a nonlegal-state border and shift the market response toward higher revenue among incumbent stores, and, separately, the composition of who enters is not something a market-size argument alone predicts.

\section{Conclusion}
\label{conclusion}

This paper studies how U.S. cannabis retail markets respond to a spatially concentrated demand shock -- captive cross-border demand from consumers who live near, but cannot legally buy from, a neighboring nonlegal state -- and asks not only how large the response is but who supplies it. Comparing New Mexico ZIP codes near a nonlegal-state line to ZIP codes near an already-legal-state line, a comparison chosen specifically to isolate captive demand from generic border proximity, the paper establishes three facts in sequence. First, captive cross-border demand exists: nonlegal-border ZIP codes generate substantially more cannabis revenue and support substantially more dispensaries, relative to population, than always-legal-border ZIP codes. Second, this additional demand is absorbed primarily through entry rather than through higher revenue at incumbent stores, consistent with a free-entry market. Third, the additional entry is not drawn uniformly from the pool of potential entrants, but the way it differs is a matter of timing rather than a permanent split: chain operators capture most of their eventual response within the first quarter after legalization, independent, single-location operators take roughly a year to catch up, and from then on the two respond identically. A simple diffusion model confirms that this is a difference in how much of each type's response was already in motion when the market opened, not in how fast either type moves once it starts.

To organize and explain this last fact, the paper builds a ZIP-code-level entry-threshold game with two ex ante heterogeneous pools of potential entrants, estimates it by simulated minimum distance at both the ZIP-code and town level, and confirms that the independent pool's estimated demand response exceeds the chain pool's in both. It then derives an explicit candidate mechanism -- established operators partly internalize the effect of a new outlet on their own nearby stores, while first-location independents do not -- and tests this mechanism's own parameter directly at the operator level. The direct test does not confirm the mechanism. The paper reports this null result in full, alongside a discussion of what it does and does not rule out, and identifies three specific, economically motivated alternative channels the available data cannot distinguish from the internalization story: financing constraints, local information advantages, and differential risk tolerance toward a policy-contingent demand source.

Using Colorado, where entry into the legal cannabis market is more restricted, the paper then compares how the market-level response to captive demand appears under a different regulatory regime. The results suggest that in Colorado, the response appears primarily through higher revenue per dispensary rather than through a larger number of dispensaries in exposed markets, interpreted as suggestive evidence that entry restrictions can limit the extent to which markets exposed to captive demand expand.

Taken together, the results show that the economic response to a spatially concentrated demand shock depends on more than its size: it depends on how freely firms can enter, and on which kind of firm -- independent or chain -- responds to the opportunity. The paper establishes both margins directly, offers one candidate explanation for the second, and is explicit that this explanation remains a hypothesis the data do not confirm. The findings also provide policymakers with information that can be used to weigh the revenue associated with market expansion, and its distribution across organizational forms, against the social costs of expanding sin-good markets.

\section{Limitations}
\label{sec:limitations}

\begin{itemize}
\item \textbf{The own-network mechanism is not empirically confirmed.} Section~\ref{sec:mechreg}'s operator-level interaction $\hat\varphi$ is statistically indistinguishable from zero in the preferred specification and significantly the wrong sign in the robustness check that includes the largest metro market. The theoretical mechanism (Section~\ref{sec:cannibalization}) remains in the paper because it is independently informative and not ruled out by an absence of network overlap, but $\hat\delta_I>\hat\delta_C$ should be read as \emph{consistent with}, not \emph{explained by}, the mechanism. $\rho$ is not estimated as a structural parameter anywhere in this paper.
\item \textbf{Neither ZIP-code exposure nor an operator's existing network is randomly assigned.} A ZIP code's nonlegal-border, always-legal-border, or interior status is fixed geography; the controls in equation~\eqref{baseline} (population, calendar-quarter, and market-age fixed effects) address observable confounds but cannot address unobserved differences correlated with which state line a ZIP code happens to sit near. Likewise, an operator's own store locations, used in Section~\ref{sec:mechreg}, are themselves the result of past entry decisions, not an experiment.
\item \textbf{The ZIP-code entry game does not have a uniquely converged equilibrium at the estimated parameters.} Section~\ref{sec:modelfit} reports this directly rather than glossing over it: multiple starting points for the fixed-point solver converge to different equilibria at the ZIP-code estimates, unlike the town-level estimation. The counterfactual in Section~\ref{sec:legalization_cf} is therefore computed from the town-level model specifically, and the ZIP-code estimation is used for its point estimates and sign, not for counterfactual simulation.
\item \textbf{$F_I, F_C$ are not separately identified from $\kappa_I,\kappa_C$.} The model speaks to net profitability differences, not to a decomposition into higher fixed costs versus higher baseline revenue for either type. A license-fee or entry-cost shifter that varies within type, not currently in this project's data, would be needed to separate them.
\item \textbf{The entry game itself has no continuation-value model, and the diffusion model of Section~\ref{sec:diffusion} is reduced-form, not structural.} The number of observed exits, spread across the full market-period-type cell count, does not support a dynamic Bayesian--Nash game with beliefs evolving over the transition; the entry-threshold game is a repeated static game describing the steady state a market settles into, not the path to it. Section~\ref{sec:diffusion} characterizes that path with a two-parameter diffusion curve fit directly to the dynamic composition coefficients, which is transparent and tightly fitting but does not model potential entrants' strategic interaction, or their beliefs about rivals, while the market is still adjusting. A fully dynamic entry game -- in which each type's entry probability depends on the current, not yet converged, count of the other type, and on a continuation value for waiting -- is a natural extension this paper's data are too short an event window to discipline.
\item \textbf{No price, quantity, or marginal-cost data.} The model cannot speak to consumer welfare, markups, or the incidence of cross-border cannabis revenue; every outcome in this paper is a revenue or count outcome.
\item \textbf{Small cross-section.} With 81--97 ZIP codes (78 towns), the model deliberately restricts its parameter count (Section~\ref{sec:identification}); a richer specification is not ruled out on economic grounds, only on the grounds that these data cannot support estimating it.
\end{itemize}

\newpage

\appendix
\section{Appendix}
\renewcommand{\thetable}{A.\arabic{table}}
\renewcommand{\thefigure}{A.\arabic{figure}}
\renewcommand{\theequation}{A.\arabic{equation}}
\setcounter{table}{0}
\setcounter{figure}{0}
\setcounter{equation}{0}

\subsection{Additional Composition Robustness}
\label{app:composition}

\begin{table}[htbp]
\centering\footnotesize
\begin{tabular}{@{}lcccc@{}}
\toprule
 & \multicolumn{2}{c}{Pooled control, ABQ included} & \multicolumn{2}{c}{Pooled control, ABQ excluded} \\
\cmidrule(lr){2-3}\cmidrule(lr){4-5}
Outcome (asinh) & Coef. & (SE) & Coef. & (SE) \\
\midrule
N independent brands & +0.244*** & (0.061) & +0.291*** & (0.044) \\
N chain brands & +0.194** & (0.095) & +0.265*** & (0.075) \\
\bottomrule
\end{tabular}

\caption{Composition result with Albuquerque included versus excluded from the pooled-control (legal-border $+$ interior) sample, full zero-inclusive 368-ZIP-code panel, inverse hyperbolic sine transformation. Albuquerque is interior by construction, so it does not affect the main (nonlegal- vs.\ always-legal-border) comparison.}
\label{tab:fact3abq}
\end{table}

\begin{table}[htbp]
\centering\footnotesize
\begin{tabular}{@{}lcccc@{}}
\toprule
 & \multicolumn{2}{c}{Main: nonlegal vs.\ legal-border towns} & \multicolumn{2}{c}{Pooled: legal+interior pooled} \\
\cmidrule(lr){2-3}\cmidrule(lr){4-5}
Outcome (asinh) & Coef. & (SE) & Coef. & (SE) \\
\midrule
Revenue per resident & +0.582 & (0.523) & +0.436 & (0.395) \\
Dispensary density & +0.199* & (0.114) & +0.170* & (0.091) \\
N independent brands & +0.318** & (0.148) & +0.452*** & (0.130) \\
N chain brands & +0.340 & (0.224) & +0.213 & (0.163) \\
\bottomrule
\end{tabular}

\caption{Main specification re-estimated at the town level: full zero-inclusive panel of 78 town-markets (Albuquerque excluded), inverse hyperbolic sine transformation, same controls as the ZIP-code main specification (no market-age control).}
\label{tab:townlevel}
\end{table}

\subsection{Additional Results for New Mexico}
\label{app:nm}

\begin{table}[htbp]
\centering\footnotesize
\begin{tabular}{@{}lcccccccc@{}}
\toprule
Sample & ZIPs & ZIP-qtrs & Avg.\ age & Pop. & Revenue & N stores & N indep. & N chain \\
\midrule
\multicolumn{9}{@{}l}{\textit{Panel A: Full sample (all 368 populated ZIP codes, zero-inclusive)}} \\
Full sample & 368 & 6256 & 5.4 & 5,588 & \$369,130 & 1.70 & 0.62 & 1.01 \\
Nonlegal-border & 92 & 1564 & 5.1 & 5,688 & \$543,232 & 2.32 & 0.99 & 1.24 \\
Always-legal-border & 117 & 1989 & 5.6 & 1,875 & \$63,114 & 0.36 & 0.15 & 0.20 \\
Interior & 159 & 2703 & 5.5 & 8,262 & \$493,573 & 2.34 & 0.75 & 1.46 \\
\midrule
\multicolumn{9}{@{}l}{\textit{Panel B: Conditional on activity (125 ZIP codes with recorded dispensary activity)}} \\
Full sample & 125 & 1745 & 5.4 & 17,267 & \$1,323,368 & 6.10 & 2.21 & 3.60 \\
Nonlegal-border & 35 & 493 & 5.1 & 15,703 & \$1,723,356 & 7.35 & 3.13 & 3.95 \\
Always-legal-border & 24 & 262 & 5.6 & 7,191 & \$479,134 & 2.72 & 1.13 & 1.51 \\
Interior & 66 & 990 & 5.5 & 20,713 & \$1,347,605 & 6.38 & 2.04 & 3.99 \\
\bottomrule
\end{tabular}

\caption{Summary statistics by exposure group, New Mexico ZIP-code panel, 2022Q2--2026Q2. Panel A: full ZIP-code panel (all 368 populated ZIP codes, zero-inclusive); Panel B: secondary conditional-on-activity sample (125 ZIP codes with at least one recorded dispensary). Revenue and store counts are ZIP-quarter means. Average dispensary age (in quarters since a storefront's own first active quarter) is computed over active dispensary-quarters only, so it is identical across the two panels: ZIP-quarters with no dispensary contribute no dispensary-age observations either way.}
\label{tab:sumstats_nm}
\end{table}

\begin{table}[htbp]
\centering
\caption{New Mexico dispensary-level border effect, noise-minimizing specification (grid-search selected)}
\label{tab:nm_dispensary_border_best}
\footnotesize
\begin{tabular}{lccccc}
\toprule
Outcome & Coef. (\%) & SE & $p$ (cl.) & $n$ & Brand FE \\
\midrule
Revenue per capita & $+0.149$ ($+16.1\%$) & 0.107 & 0.162 & 10725 & Yes \\
Revenue (per dispensary) & $+0.150$ ($+16.2\%$) & 0.103 & 0.144 & 10739 & Yes \\
Local dispensary density & $-0.154$ ($-14.3\%$) & 0.160 & 0.336 & 10739 & No \\
\bottomrule
\end{tabular}
\begin{minipage}{0.98\textwidth}
\vspace{0.3em}
\footnotesize
\textit{Notes:} The sample consists of geocoded New Mexico dispensaries with valid coordinates (1,192 of 1,253). Border =
dispensary within 75 miles of the nonlegal-state line, with $\geq$1,000 nonlegal-state residents within that
distance. All specifications include a population control in logs and quarter fixed effects and use county-clustered standard errors;
chain/brand fixed effects are included for the two revenue outcomes only (see text).
\end{minipage}
\end{table}

\begin{table}[htbp]
\centering
\caption{Dose-response of dispensary outcomes to distance from the nonlegal-state line, New Mexico}
\label{tab:nm_doseresponse}
\footnotesize
\begin{tabular}{lccc}
\toprule
Distance band (mi) & Revenue per capita & Revenue per dispensary & Dispensary density \\
\midrule
0-25 & ref. & ref. & ref. \\
25-50 & $+0.275$ ($+31.6\%$), $p=0.244$ & $+0.277$ ($+31.9\%$), $p=0.235$ & $+0.083$ ($+8.6\%$), $p=0.827$ \\
50-75 & $+0.054$ ($+5.5\%$), $p=0.810$ & $+0.044$ ($+4.5\%$), $p=0.838$ & $-0.217$ ($-19.5\%$), $p=0.276$ \\
75-100 & $+0.166$ ($+18.0\%$), $p=0.432$ & $+0.150$ ($+16.2\%$), $p=0.476$ & $-0.634$ ($-46.9\%$), $p=0.002$ \\
100-125 & $+0.289$ ($+33.5\%$), $p=0.315$ & $+0.291$ ($+33.8\%$), $p=0.308$ & $-0.473$ ($-37.7\%$), $p=0.095$ \\
125-150 & $+0.257$ ($+29.3\%$), $p=0.244$ & $+0.259$ ($+29.6\%$), $p=0.233$ & $-0.536$ ($-41.5\%$), $p=0.002$ \\
150-175 & $+0.116$ ($+12.3\%$), $p=0.561$ & $+0.109$ ($+11.5\%$), $p=0.580$ & $-0.675$ ($-49.1\%$), $p<0.001$ \\
\bottomrule
\end{tabular}
\begin{minipage}{0.98\textwidth}
\vspace{0.3em}
\footnotesize
\textit{Notes:} Each column reports a separate regression of the indicated log outcome on the full set of distance-band
indicators, with 0--25 miles as the omitted reference category, a log population control, and quarter fixed effects.
The 175--200-mile band is excluded from the sample (see text). Standard errors are clustered at the county level.
\end{minipage}
\end{table}

\begin{table}[htbp]
\centering
\caption{Dose-response of the border premium to nearby nonlegal-state population, New Mexico}
\label{tab:nm_nonlegalpop_doseresponse}
\footnotesize
\begin{tabular}{lcccc}
\toprule
 & \multicolumn{2}{c}{Revenue per capita} & \multicolumn{2}{c}{Dispensary density} \\
\cmidrule(lr){2-3}\cmidrule(lr){4-5}
Nonlegal pop.\ within 75 mi (000s) & Coef. & SE & Coef. & SE \\
\midrule
0 (reference) & ref. & & ref. & \\
0--50 & $-0.019$ & (0.221) & $+0.187$ & (0.138) \\
50--250 & $+0.869^{**}$ & (0.356) & $+0.465^{***}$ & (0.179) \\
250+ & $+0.859^{***}$ & (0.238) & $+0.437^{***}$ & (0.150) \\
\midrule
Quarter fixed effects & Yes & & Yes & \\
Observations & 1,704 & & 1,704 & \\
Number of ZIP codes & 116 & & 116 & \\
\bottomrule
\end{tabular}
\begin{minipage}{0.92\textwidth}
\vspace{0.4em}\footnotesize
\textit{Notes:} Each column reports $\beta_k$ from $\log(y_{zt}) = \alpha + \sum_k \beta_k\,\mathbbm{1}[\text{bin}_z=k] + \gamma\log(\text{Population}_z) + \delta_t + \varepsilon_{zt}$, estimated on the 116 New Mexico ZIP codes with any dispensary activity. $\text{Population}_z$ is the ZIP's own resident population. Quarter fixed effects, county-clustered standard errors. $^{*}p<0.10$, $^{**}p<0.05$, $^{***}p<0.01$.
\end{minipage}
\end{table}


\begin{figure}[htbp]
\centering
\includegraphics[width=\textwidth]{fig_nm_zip_noncannabis_placebo_legaleventstudy.png}
\caption{Border coefficient by calendar quarter, non-cannabis sectors, New Mexico ZIP-quarter panel.
$\text{Border}_z$ is interacted with each calendar quarter in the sample, anchored to New Mexico's own
legalization date; 95\% confidence intervals are ZIP-clustered.}
\label{fig:noncannabis_eventstudy}
\end{figure}

\begin{figure}[htbp]
\centering
\includegraphics[width=\textwidth]{fig_nm_zip_eventstudy_50_100mi.png}
\caption{Dynamic border premium under the 50-mile and 100-mile border definitions. Reference quarter: 2022Q2; county-clustered 95\% confidence intervals.}
\label{fig:eventstudy_50_100mi}
\end{figure}

\begin{figure}[htbp]
\centering
\includegraphics[width=\textwidth]{fig_nm_doseresponse.png}
\caption{Dose-response of dispensary outcomes to distance from the nonlegal-state line, New Mexico. The omitted reference category is 0--25 miles. County-clustered 95\% confidence intervals.}
\label{fig:nm_doseresponse}
\end{figure}

\begin{figure}[htbp]
\centering
\includegraphics[width=\textwidth]{fig_nm_nonlegalpop_doseresponse.png}
\caption{Dose-response of the New Mexico border premium to nearby nonlegal-state population. New Mexico ZIP codes are grouped by the number of Texas, Oklahoma, or Utah residents living within 75 miles of the ZIP centroid; the reference category consists of ZIP codes with essentially no such residents. Quarter fixed effects, county-clustered 95\% confidence intervals.}
\label{fig:nm_nonlegalpop_doseresponse}
\end{figure}

\subsection{Cross-State Entry-Regime Comparisons}
\label{sec:crossstate}

\begin{figure}[htbp]
\centering
\includegraphics[width=\textwidth]{fig1_southwest_map.png}
\caption{Dispensary density and the nonlegal-state border, New Mexico and
Colorado, 2025. Shading indicates dispensaries per 1,000 residents in 2025. Red outlines identify counties containing at least one
border-exposed dispensary (19 in New Mexico and 17 in Colorado).}
\label{fig:setting}
\label{map}
\end{figure}

\begin{table}[htbp]
\centering
\caption{Cross-state entry-regime comparison: county-level sample, New Mexico and Colorado}
\label{tab:sumstats_state}
\footnotesize
\begin{tabular}{lcc}
\toprule
 & New Mexico & Colorado \\
\midrule
Entry regime & Free entry & Local option / capped \\
First recreational retail sales & 1 April 2022 & 1 January 2014 \\
\addlinespace
Counties (total) & 33 & 64 \\
\quad Border counties (geography) & 11 & 23 \\
\addlinespace
Counties in estimation sample & 32 & 42 \\
County--quarters & 360 & 530 \\
\quad Border county--quarters & 118 & 149 \\
\addlinespace
\multicolumn{3}{l}{\textit{County-quarter means}} \\
\quad Revenue per capita (\$) & 66.5 & 153.1 \\
\quad Dispensaries per 1,000 residents & 0.432 & 0.438 \\
\quad Revenue per dispensary (\$000s) & 213.2 & 524.8 \\
\quad County population & 65,955 & 125,571 \\
\quad Nonlegal pop.\ within 75 mi, border counties (000s) & 194.5 & 37.7 \\
\bottomrule
\end{tabular}
\begin{minipage}{0.86\textwidth}
\vspace{0.4em}\footnotesize
\textit{Notes:} County-quarter panel, 2022Q2 through 2025Q2, the sample used for the cross-state entry-regime comparison. Colorado's own market opened in 2014, but it is restricted to 2022Q2 onward.New Mexico is coded free entry; Colorado is coded locally capped because it delegates licensing to local governments and, as of 2019, 48.4 percent of its counties banned recreational cannabis businesses outright \citep{PayanBrownSong2021}.
\end{minipage}
\end{table}

\begin{table}[htbp]
\centering
\caption{Border $\times$ Restricted interaction: New Mexico vs. Colorado}
\label{tab:regime_interaction}
\footnotesize
\begin{tabular}{lc}
\toprule
Outcome & NM+CO \\
\midrule
Revenue per capita & $-0.458$ \\
 & (0.388) \\
Dispensary density & $-0.932^{*}$ \\
 & (0.488) \\
Revenue per store & $+0.465$ \\
 & (0.314) \\
\bottomrule
\end{tabular}
\begin{minipage}{0.98\textwidth}
\vspace{0.5em}

\fontsize{9pt}{11pt}\selectfont
\textit{Notes:} Each row is the interaction coefficient
$\theta$ from a separate regression for the named outcome; the number in
parentheses directly below it is the standard error, clustered at the county
level. Restricted$_s$ = 1 for Colorado and 0 for New Mexico. Both Border$_c$ and Restricted$_s$ are
time-invariant in this sample. $\theta$ is therefore interpreted as the
difference in the estimated border premium between Colorado and New Mexico,
rather than as an unqualified causal treatment effect.
$^{*}p<0.10$, $^{**}p<0.05$, $^{***}p<0.01$.
\end{minipage}
\end{table}

\begin{table}[htbp]
\centering
\caption{Model-implied dispensary-revenue benchmark, Colorado border counties, 2024}
\label{tab:counterfactual}
\footnotesize
\begin{tabular}{lc}
\toprule
Quantity & Value \\
\midrule
Dispensary-density interaction coefficient $\theta$ & $-0.932$ \\
Free-entry benchmark multiplier $e^{-\theta}$ & 2.54 \\
\midrule
Actual 2024 revenue, CO border counties & \$324.2 M \\
Free-entry-benchmark counterfactual revenue (NM) & \$823.0 M \\
Model-implied additional dispensary revenue & +\$498.8 M \\
95\% county-cluster bootstrap interval & (\$-1.5 M, +\$2{,}526.0 M) \\
\bottomrule
\end{tabular}
\begin{minipage}{0.98\textwidth}
\vspace{0.5em}

\fontsize{9pt}{11pt}\selectfont
\textit{Notes:} The benchmark uses the \emph{dispensary-density} interaction
($\theta=-0.932$, Table~\ref{tab:regime_interaction}), the outcome for which the
difference across entry regimes is statistically significant. It asks how much
more dispensary revenue Colorado's border counties would have generated in 2024
if their density border premium had matched New Mexico's rather than the
restricted-entry premium, holding revenue per store at its observed Colorado
level. The multiplier $e^{-\theta}$ is applied to each Colorado border county's
own 2024 dispensary revenue and summed; New Mexico is the free-entry benchmark
state and contributes only to the estimation of $\theta$. Colorado is restricted
to 2022Q2 onward (Section~\ref{sec:crossstate}). The bootstrap interval resamples
counties with replacement within each of the two states (500 replicates),
refitting the density regression each time; it is wide and reaches down to zero
because $\theta$ is only marginally significant ($p\approx0.06$). The dollar
figure should therefore be read as an illustrative benchmark rather than an
identified estimate of foregone revenue.
\end{minipage}
\end{table}


\subsection{Technical Appendix: Estimation of the Entry Game}
\label{app:computation}

\subsubsection{Vectorized Newton solution of the Bayesian--Nash fixed point}
For a fixed $\theta$, every market-period's fixed point $(N^{I*}_{zt},N^{C*}_{zt})$ solves its own two-equation system, independent of every other market-period's, so the $n$-market-period problem decomposes into $n$ independent $2\times2$ systems, solved in one vectorized pass. Writing $F_1 = N^I - R_I(N^C)$ and $F_2 = N^C - R_C(N^I)$, the Jacobian of $(F_1,F_2)$ with respect to $(N^I,N^C)$ is
\[
J = \begin{pmatrix} 1 & -R_I'(N^C) \\ -R_C'(N^I) & 1 \end{pmatrix}, \qquad R_I'(N^C) = -\gamma_{IC}\sum_{k=1}^{\bar K}\Lambda\big(\pi^I(k)\big)\big(1-\Lambda\big(\pi^I(k)\big)\big),
\]
and symmetrically for $R_C'(N^I)$. With $a \equiv -R_I'(N^C)$, $b \equiv -R_C'(N^I)$, the Newton step has the closed form
\[
\Delta N^I = \frac{-F_1 + aF_2}{1-ab}, \qquad \Delta N^C = \frac{bF_1 - F_2}{1-ab},
\]
applied elementwise across all market-periods simultaneously. The solver runs six damped best-response iterations, to enter Newton's basin of attraction, followed by up to 60 Newton steps with a bounded step size, clipping $(N^I,N^C)$ to $[0,\bar K]$ at each iterate; $\bar K=80$.

\subsubsection{Sparse two-way fixed-effects PPML}
\label{app:ppml}
The mechanism regression's risk-set panel (Section~\ref{sec:mechreg}) has several hundred thousand rows and, with operator and quarter fixed effects, hundreds of parameters -- large enough that dense dummy-variable encoding is not memory-feasible. The estimator implements Poisson IRLS directly on a sparse design matrix: at each iterate, the weighted normal equations $X'WX\,\Delta\beta = X'W(z-X\beta)$ ($W=\text{diag}(\hat\mu)$, the canonical Poisson IRLS weight; $z$ the working response) are formed via sparse matrix products, so only the resulting small dense system is ever solved directly. Standard errors are cluster-robust by operator, computed with the same sparse machinery.

\newpage
\begin{singlespace}
\bibliographystyle{aer}

\bibliography{bibliography.bib}

\end{singlespace}

\end{document}

\section*{Declarations}
\begin{singlespace}
\noindent \textbf{Data availability.} This study uses only publicly available data. New Mexico cannabis sales and licensed-establishment records are published by the New Mexico Cannabis Control Division; Colorado county-level marijuana sales and licensee records are published by the Colorado Department of Revenue; population, geographic, and business-pattern data are obtained from the U.S. Census Bureau and the SimpleMaps U.S. ZIP Code Database. The code used to construct the analysis panels and reproduce all tables and figures is available from the author upon request.

\smallskip
\noindent \textbf{Conflicts of interest.} The author has no competing interests relevant to the content of this article.

\smallskip
\noindent \textbf{Funding.} The author received no specific funding for this work.

\smallskip
\noindent \textbf{Use of AI tools.} The author used generative AI assistance for LaTeX formatting, code implementation, and copy-editing. All analysis, results, and interpretation are the author's own.
\end{singlespace}






